% =====================================================================
%  Surrogate-Accelerated Conjunction Assessment -- running project log
%
%  Reference sheet for the write-up. Compiles standalone:
%      pdflatex project_log.tex
%  or % =====================================================================
%  Surrogate-Accelerated Conjunction Assessment -- running project log
%
%  Reference sheet for the write-up. Compiles standalone:
%      pdflatex project_log.tex
%  or % =====================================================================
%  Surrogate-Accelerated Conjunction Assessment -- running project log
%
%  Reference sheet for the write-up. Compiles standalone:
%      pdflatex project_log.tex
%  or % =====================================================================
%  Surrogate-Accelerated Conjunction Assessment -- running project log
%
%  Reference sheet for the write-up. Compiles standalone:
%      pdflatex project_log.tex
%  or \input{project_log} into a larger document.
%
%  Every table and section is \label'd for cross-referencing, using the
%  prefixes tab: and sec:. Numbers here are measured, not estimated;
%  where a number is an assumption it is marked as such.
% =====================================================================
\documentclass[11pt]{article}

\usepackage[margin=1in]{geometry}
\usepackage{amsmath,amssymb}
\usepackage{booktabs}
\usepackage{longtable}
\usepackage{siunitx}
\usepackage{hyperref}
\usepackage{xcolor}

\newcommand{\Pc}{P_c}
\newcommand{\worked}[1]{\textcolor{teal}{\textbf{Worked:}} #1}
\newcommand{\failed}[1]{\textcolor{red}{\textbf{Did not work:}} #1}
\newcommand{\gotcha}[1]{\textcolor{orange}{\textbf{Gotcha:}} #1}

\title{Surrogate-Accelerated Conjunction Assessment\\
       \large Running project log and results reference}
\author{Anurag Paudel}
\date{Last updated: 8 September 2026}

\begin{document}
\maketitle

\begin{abstract}
Working log for a project that estimates satellite collision probability
$\Pc$ using a Gaussian-process surrogate trained at space-filling design
points, benchmarked against a brute-force Monte Carlo baseline. This
document records what was built, what was measured, what failed, and which
claims are verified versus assumed. It is a reference sheet for the
write-up, not the write-up itself.
\end{abstract}

\tableofcontents

% =====================================================================
\section{Summary of headline results}\label{sec:headline}

\begin{longtable}{p{0.42\textwidth} p{0.5\textwidth}}
\toprule
\textbf{Result} & \textbf{Value} \\
\midrule
\endhead
TLE uncertainty, calibrated from data &
  $\sigma_R = \SI{95}{\meter}$, $\sigma_T = \SI{953}{\meter}$,
  $\sigma_N = \SI{143}{\meter}$ \\
Monte Carlo cost, $\Pc \approx 4\times10^{-6}$ &
  $\sim$\num{22000000} draws for \SI{10}{\percent} relative error;
  $\sim2.2\times10^{9}$ for \SI{1}{\percent} \\
Monte Carlo at \num{100000} draws & \textbf{zero hits} -- no estimate at all \\
Dynamic range of $\Pc$ over the parameter box & \num{12.8} orders of magnitude \\
Fraction of design points brute-force MC can label at $10^8$ draws &
  \SI{28}{\percent} \\
Dimensional reduction & $6 \to 4$ parameters, exact to $\sim10^{-12}$ \\
Design points saved by the reduction & $8$--$14\times$ at equal fill distance \\
Parallel tempering vs.\ best-of-20 random LHD & $2.1$--$3.1\times$ better $\psi$ \\
\bottomrule
\end{longtable}

% =====================================================================
\section{Phase 0: environment}\label{sec:phase0}

Python \num{3.9.6} (macOS system interpreter), virtual environment in
\texttt{venv/}. Dependencies: \texttt{sgp4}, \texttt{numpy},
\texttt{pandas}, \texttt{matplotlib}, \texttt{requests}, later
\texttt{scipy} and \texttt{pybind11}.

\gotcha{Python 3.9 is end-of-life, so \texttt{pip} resolved older builds
(\texttt{numpy 2.0.2}, \texttt{matplotlib 3.9.4}). The system \texttt{ssl}
module is compiled against LibreSSL 2.8.3, so every network call emits a
\texttt{NotOpenSSLWarning}. Harmless but noisy; a rebuild on Python
3.11+ would remove both issues.}

% =====================================================================
\section{Phase 1: data ingestion and propagation}\label{sec:phase1}

\subsection{TLE format}

Two-line element sets are fixed-column records, 69 characters per line.
Fields are located by column position, not delimiter. The six orbital
elements live on line 2: inclination, RAAN, eccentricity, argument of
perigee, mean anomaly, and mean motion (which stands in for the
semi-major axis, recovered as $a = (\mu/n^2)^{1/3}$).

Four format traps, all of which the parser handles explicitly:
\begin{enumerate}
  \item \textbf{Implied decimal points.} Eccentricity \texttt{0006703}
        means \num{0.0006703}. The $B^*$ and second-derivative fields use
        assumed-decimal exponential notation: \texttt{30074-3} is
        $0.30074\times10^{-3}$.
  \item \textbf{Two-digit epoch year.} 57--99 map to 19xx, 00--56 to 20xx.
  \item \textbf{Elements are Brouwer mean elements}, defined by the SGP4
        theory itself. Propagating them with a two-body Kepler propagator
        produces kilometres of error immediately. The TLE and SGP4 are a
        matched pair.
  \item \textbf{Output frame is TEME}, not J2000 and not ECEF.
\end{enumerate}

\subsection{What worked}

\worked{The ISS validation passed every check. Propagating three days
forward from epoch at one-minute cadence gave a mean altitude of
\SI{419.4}{\kilo\meter} (range \num{409.3}--\SI{427.7}{\kilo\meter}),
speed \num{7.645}--\SI{7.666}{\kilo\meter\per\second}, and a period of
\SI{92.96}{\minute} from the elements. An \emph{independent} empirical
period, measured by counting geocentric-radius minima over the
propagation, gave \SI{92.91}{\minute} -- agreement to
\SI{0.05}{\minute}, confirming the propagator is not simply echoing its
input.}

\subsection{Problems encountered}

\failed{My hand-typed ISS test fixture failed its own checksum on both
lines. The checksum implementation was correct -- it validated against
live CelesTrak data on both lines -- so the fixture's two check digits
were recomputed. Worth recording as a methodological point: when a test
fails, verify the test before ``fixing'' the code.}

\gotcha{CelesTrak returns HTTP 200 with a plaintext message, not an error
status, when a query matches nothing or the client is throttled. The
fetcher content-checks the response rather than trusting the status code.}

\gotcha{Space-Track registration completed \emph{instantly}, contrary to
documentation describing manual review taking days. Worth noting because
project planning had treated it as the long-lead dependency.}

% =====================================================================
\section{Phase 2: Monte Carlo baseline}\label{sec:phase2}

\subsection{What the public data actually contains}

The \texttt{cdm\_public} feed has \textbf{16 fields and no covariance
matrices}. It is a screening summary, not a full CCSDS Conjunction Data
Message: no state vectors, no epochs, no RTN covariance. Available:
\texttt{CDM\_ID}, \texttt{TCA}, \texttt{PC}, \texttt{MIN\_RNG},
\texttt{EMERGENCY\_REPORTABLE}, object IDs, names, types, RCS size class,
and exclusion volumes.

Three consequences shaped everything downstream:
\begin{enumerate}
  \item Covariances must be \textbf{synthesized}. This is a stated
        modelling assumption, not a data product.
  \item \texttt{PC} is the 18th Space Defense Squadron's own published
        probability. It is a reference point, \emph{not} a validation
        target, because it was computed from the covariance we do not
        have.
  \item The feed is pre-filtered to high-risk events: every sampled row
        had \texttt{EMERGENCY\_REPORTABLE = Y} and $\texttt{PC} > 10^{-4}$.
        There are \textbf{no negative examples}, which constrains the
        Phase 4 classifier.
\end{enumerate}

\gotcha{\texttt{EXCL\_VOL} is a kilometre-scale \emph{screening} volume
keyed to object class (debris 1, rocket body 3, payload 5), \textbf{not} a
hard-body radius. Since $\Pc \propto \text{HBR}^2$, using it as a hard-body
radius would have inflated every probability by roughly six orders of
magnitude. Values are fully determined by object type, which is what gave
it away.}

\subsection{Rebuilding real conjunctions from TLEs}

Both objects of each event were propagated to TCA with SGP4 and the
separation compared against the reported \texttt{MIN\_RNG}
(Table~\ref{tab:rebuild}).

\begin{table}[htbp]
\centering
\caption{Rebuilt versus reported miss distance, 37 deduplicated events.}
\label{tab:rebuild}
\begin{tabular}{lr}
\toprule
Quantity & Value (\si{\kilo\meter}) \\
\midrule
Reported miss distance (\texttt{MIN\_RNG}) range & 0.008--4.510 \\
Median $|$rebuilt $-$ reported$|$ & 0.936 \\
Mean $|$rebuilt $-$ reported$|$ & 2.219 \\
Range of $|$rebuilt $-$ reported$|$ & 0.046--20.798 \\
\bottomrule
\end{tabular}
\end{table}

This gap is the central motivating result: \textbf{reported miss distances
are hundreds of metres, while TLE-based rebuilds disagree by kilometres}.
Collision assessment cannot be done by propagating TLEs and measuring a
distance, because the error is an order of magnitude larger than the
quantity being measured. That is why the problem is probabilistic.

\subsection{Two observations worth keeping}

\textbf{Every rebuild error was positive.} The rebuilt miss is always
\emph{larger} than reported, never smaller. This is a genuine statistical
effect rather than a bias in the code: when the true miss is near zero,
adding random error to either object almost always increases the
separation, because distance is a norm. The Monte Carlo reproduces the
same asymmetry.

\textbf{Events are filed multiple times.} Each conjunction appears twice
(once with each object as primary) and is re-filed as the estimate is
refined, each revision at a slightly different TCA. Deduplication
therefore groups by unordered object pair and merges revisions within a
\SI{15}{\minute} window. An exact-TCA key silently missed the revisions.

\subsection{Covariance model (Option B)}

Diagonal in each object's RTN frame,
\begin{equation}
  C_{\text{RTN}} = \operatorname{diag}\!\left(\sigma_R^2,\;
    (k_T\sigma_R)^2,\; (k_N\sigma_R)^2\right),
  \label{eq:covb}
\end{equation}
with anisotropy ratios \emph{fixed} at $k_T = 10$, $k_N = 1.5$ and only
the scale $\sigma_R$ fitted. Twelve free parameters cannot be identified
from one scalar observable per event; one can. In-track error dominates
because a small period error integrates into a growing along-track lag.

Each object's covariance is rotated into the inertial frame \emph{before}
summing. The two objects occupy different orbital planes, so adding them
componentwise in RTN would be wrong.

\begin{table}[htbp]
\centering
\caption{Calibrated covariance scale, 37 events, ratios held fixed.}
\label{tab:calib}
\begin{tabular}{lS[table-format=4.0]}
\toprule
Component & {Value (\si{\meter})} \\
\midrule
$\sigma_R$ (radial)      & 95 \\
$\sigma_T$ (in-track)    & 953 \\
$\sigma_N$ (cross-track) & 143 \\
\bottomrule
\end{tabular}
\end{table}

\worked{These magnitudes are consistent with published TLE accuracy
($\sim$\SI{1}{\kilo\meter} in-track at a few days) and were derived from
the project's own measurements, not assumed.}

\failed{Per-event correlation between predicted and observed error is
$-0.53$ -- \emph{negative}. The model captures how large TLE error is in
aggregate but gets the per-event ranking backwards. The honest claim is
that it reproduces the scale of uncertainty, not which conjunctions are
worst. Cause: a single scalar observable per event is dominated by the
reported miss magnitude, and the model's per-event variation keys off the
same quantity in the opposite direction.}

\failed{Residual-versus-TLE-age correlation is only $-0.16$, so the data
does \emph{not} support an explicit time-growth term. Staying with the
fixed-scale model is an empirical decision, not a simplification.}

\subsection{Monte Carlo and its cost}

The short-term encounter model applies: at
\SIrange{10}{15}{\kilo\meter\per\second} relative speed the encounter
lasts milliseconds, relative motion is effectively rectilinear, and the
three-dimensional problem collapses onto the two-dimensional encounter
plane perpendicular to the relative velocity. $\Pc$ is then the integral
of the projected Gaussian over a disk of the combined hard-body radius.

Three independent estimators were implemented: batched Monte Carlo, polar
quadrature, and a small-disk closed form valid when
$\text{HBR} \ll \sigma$.

\worked{Monte Carlo and quadrature agree within Monte Carlo error bars on
every real event, maximum $|z| = 1.4$. Two methods sharing no code path,
so the encounter geometry and the sampler are cross-validated.}

\begin{table}[htbp]
\centering
\caption{Brute-force Monte Carlo convergence at $\Pc = 4.5\times10^{-6}$.
         This is the cost the surrogate replaces.}
\label{tab:mccost}
\begin{tabular}{rlS[table-format=3.1]}
\toprule
{Draws} & {$\Pc$ estimate} & {Relative error (\si{\percent})} \\
\midrule
\num{10000}     & \num{0.0}          & 100.0 \\
\num{100000}    & \num{0.0}          & 100.0 \\
\num{1000000}   & \num{6.0e-6}       & 33.1 \\
\num{10000000}  & \num{2.7e-6}       & 40.1 \\
\num{50000000}  & \num{4.56e-6}      & 1.2 \\
\bottomrule
\end{tabular}
\end{table}

At \num{100000} draws the estimator returns \textbf{zero hits} -- not a
poor answer, but no answer. Reaching \SI{10}{\percent} relative error
requires $\sim$\num{22000000} draws and \SI{1}{\percent} requires
$\sim2.2\times10^{9}$, per conjunction, against a catalogue of thousands
of screened pairs.

\subsection{Known caveats}\label{sec:caveats}

\begin{enumerate}
  \item Computed $\Pc$ runs $10$--$100\times$ \emph{below} the published
        \texttt{PC}. The dominant driver is almost certainly the hard-body
        radius: RCS-derived values of \SIrange{1}{4}{\meter} are likely
        below operational values, and $\Pc \propto \text{HBR}^2$, so a
        $10\times$ error in HBR is $100\times$ in $\Pc$. Agreement with
        \texttt{PC} must \emph{not} be presented as validation.
  \item Per-event calibration correlation is negative (above).
  \item One event showed a \SI{20.8}{\kilo\meter} rebuild error, far
        outside plausible TLE error. Likely a manoeuvre or stale elements;
        should be identified before it contaminates Phase 4.
  \item Slow encounters (relative velocity below
        \SI{1}{\kilo\meter\per\second}) violate the short-term encounter
        assumption and are currently filtered out.
\end{enumerate}

% =====================================================================
\section{Phase 3: space-filling designs and the surrogate}\label{sec:phase3}

\subsection{A negative result that redirected the phase}\label{sec:negative}

The intuitive reading of ``use a space-filling design instead of Monte
Carlo'' does \textbf{not} work. $\Pc$ is a rare-event indicator integral
over a microscopic disk. A 256-point design returns zero hits just as
\num{100000} random draws do, and in fact more reliably, because a
space-filling design deliberately spreads points \emph{away} from any
small region. Quasi-Monte Carlo does not rescue this either: its
advantage requires a smooth integrand, and an indicator function over a
tiny region is the opposite.

The construction that does work is a surrogate over the
\textbf{encounter-parameter space}, trained on expensive evaluations at
design points and amortised across the catalogue. This is recorded because
it is the kind of dead end worth reporting in a write-up rather than
quietly discarding.

\subsection{Dimensional reduction, verified rather than assumed}

$\Pc$ depends on six raw parameters: the miss vector $\mu$ (2), the
projected covariance $C$ (3), and the hard-body radius $R$ (1). Only
\textbf{four} are independent, because the defining integral
\begin{equation}
  \Pc = \int_{\|x\| \le R} \mathcal{N}(x;\, \mu,\, C)\, \mathrm{d}x
  \label{eq:pcdef}
\end{equation}
has a rotation-symmetric domain and is covariant under uniform scaling of
all lengths. Rotating into the eigenbasis of $C$ removes one parameter;
dividing every length by $R$ removes another, leaving
$(\mu_1/R,\, \mu_2/R,\, \sigma_1/R,\, \sigma_2/R)$.

This is an exact identity, and it was verified numerically rather than
argued from algebra (Table~\ref{tab:invariance}).

\begin{table}[htbp]
\centering
\caption{Numerical verification of the invariances underlying the
         $6 \to 4$ reduction. Errors are quadrature noise, not model
         error.}
\label{tab:invariance}
\begin{tabular}{lrS[table-format=1.1e-2]}
\toprule
Claim & {Cases} & {Max.\ relative error} \\
\midrule
$\Pc$ invariant under rotation of $\mu$ and $C$ & 300 & 8.3e-13 \\
$\Pc$ invariant under uniform length scaling    & 300 & 9.8e-13 \\
4-D reduction reproduces $\Pc$                  & 200 & 8.4e-13 \\
$\Pc$ even in each $\mu$ component (eigenbasis) & 200 & 6.2e-16 \\
$6 \to 4 \to \Pc$ round trip                    & 200 & 5.5e-13 \\
\bottomrule
\end{tabular}
\end{table}

\worked{The reduction is exact to machine precision, so it costs no
accuracy and saves substantial cost. Measured fill distance at $k=4$,
$n=64$ is \num{0.397}; matching that in six dimensions requires
$\sim$\numrange{500}{900} points, i.e.\ $8$--$14\times$ more expensive
evaluations for identical answers.}

The identity holds \emph{given the model}: circular hard body, Gaussian
uncertainty, short-term encounter. Slow encounters fall outside it. The
six-dimensional space is retained as an ablation so the saving can be
demonstrated empirically rather than asserted.

\subsection{The 4-D parameter space}

\begin{table}[htbp]
\centering
\caption{Primary parameter space. Ranges from 40 rebuilt conjunctions,
         widened. Anisotropy is parameterised as a ratio so
         $\sigma_2 \le \sigma_1$ holds by construction.}
\label{tab:paramspace}
\begin{tabular}{llll}
\toprule
Index & Parameter & Meaning & Range \\
\midrule
0 & $\log_{10}(\sigma_1/R)$        & uncertainty vs.\ object size & 32 -- 3162 \\
1 & $\log_{10}(\sigma_2/\sigma_1)$ & anisotropy                   & 0.032 -- 1 \\
2 & $\mu_1/\sigma_1$               & miss along major axis        & 0 -- 6 \\
3 & $\mu_2/\sigma_2$               & miss along minor axis        & 0 -- 6 \\
\bottomrule
\end{tabular}
\end{table}

\subsection{Design generation}

The RUSIS parallel-tempering MaxPro code was reused directly. The
pybind11 port compiled without modification and its designs come out
already scaled to $[0,1]$ as $(\ell - 0.5)/n$.

\worked{An independent \texttt{numpy} reimplementation of the MaxPro
criterion $\psi$ agrees with the C\texttt{++} extension to six decimal
places, cross-checking the original research code.}

\begin{table}[htbp]
\centering
\caption{Design quality at $k=6$. Lower $\psi$ and $\phi_p$ are better;
         larger minimum distance is better. Random LHD is the best of 20
         draws.}
\label{tab:designquality}
\begin{tabular}{rlSSS}
\toprule
{$n$} & Design & {$\psi$ (MaxPro)} & {$\phi_p$} & {Min.\ distance} \\
\midrule
64  & MaxPro (PT) & 37.7297 & 1.7649 & 0.3694 \\
64  & Random LHD  & 81.2452 & 2.3531 & 0.2684 \\
64  & Uniform     & 508.8451 & 2.8862 & 0.2118 \\
\midrule
256 & MaxPro (PT) & 79.6778 & 2.4753 & 0.2162 \\
256 & Random LHD  & 245.0499 & 3.4067 & 0.1533 \\
256 & Uniform     & 415.3693 & 3.3181 & 0.1615 \\
\bottomrule
\end{tabular}
\end{table}

Parallel tempering beats the best of 20 random Latin hypercubes by
$2.1$--$3.1\times$ on $\psi$, at \SIrange{0.2}{1.4}{\second} per design --
a cost paid once, since a design is a static artefact.

\subsection{Why training labels come from quadrature, not Monte Carlo}
\label{sec:labels}

$\Pc$ spans \textbf{12.8 orders of magnitude} across the parameter box
($2.1\times10^{-17}$ to $1.3\times10^{-4}$ on the $n=64$ design). Brute
force cannot label it:

\begin{table}[htbp]
\centering
\caption{Fraction of the 64 design points at which brute-force Monte Carlo
         yields a usable estimate (more than 10 expected hits).}
\label{tab:mcunusable}
\begin{tabular}{rr}
\toprule
{Draws per point} & {Points labelled (of 64)} \\
\midrule
$10^6$ & 6 \phantom{0}(\SI{9}{\percent}) \\
$10^7$ & 13 (\SI{20}{\percent}) \\
$10^8$ & 18 (\SI{28}{\percent}) \\
\bottomrule
\end{tabular}
\end{table}

Nineteen of 64 points lie below $\Pc = 10^{-12}$, which no Monte Carlo
will ever reach. Labels therefore come from the polar quadrature, which is
exact here and agrees with Monte Carlo wherever Monte Carlo works at all
(Section~\ref{sec:phase2}).

This strengthens rather than weakens the motivation: brute-force Monte
Carlo is not merely \emph{expensive} for rare events, it is
\textbf{unusable} across most of the parameter space. The surrogate
produces estimates where brute force produces nothing. The quadrature
serves as exact ground truth in this Gaussian testbed; in the general case
(non-Gaussian uncertainty, long-duration encounters, non-circular bodies)
no closed form exists and the expensive sampler is the only option.

\subsection{Surrogate}

Gaussian process with an anisotropic (ARD) squared-exponential kernel,
hyperparameters by maximising the log marginal likelihood. Targets are
$\log_{10}\Pc$ rather than $\Pc$: with 13 orders of magnitude of dynamic
range, absolute error is meaningless and a GP on the raw value would be
dominated by a handful of large-$\Pc$ points. The GP is plain
\texttt{numpy} (one Cholesky factorisation, two triangular solves); only
the optimiser comes from \texttt{scipy}.

% --- RESULTS TABLE INSERTED BELOW WHEN THE BENCHMARK RUN COMPLETES ---
\input{surrogate_results}

% =====================================================================
\section{Reproducing everything}\label{sec:repro}

\begin{longtable}{p{0.42\textwidth} p{0.5\textwidth}}
\toprule
\textbf{Command} & \textbf{Produces} \\
\midrule
\endhead
\texttt{python src/validate\_iss.py}    & ISS propagation sanity checks \\
\texttt{python src/probe\_cdm.py}       & \texttt{cdm\_public} field inventory \\
\texttt{python src/rebuild\_events.py}  & Table~\ref{tab:rebuild} \\
\texttt{python src/calibrate.py}        & Table~\ref{tab:calib} \\
\texttt{python src/run\_montecarlo.py}  & Table~\ref{tab:mccost} \\
\texttt{python src/run\_surrogate.py}   & Table~\ref{tab:surrogate} \\
\texttt{python tests/test\_paramspace.py} & Table~\ref{tab:invariance} \\
\bottomrule
\end{longtable}

Test suites: \texttt{tests/test\_tle.py} (31),
\texttt{tests/test\_covariance.py} (22),
\texttt{tests/test\_montecarlo.py} (18),
\texttt{tests/test\_paramspace.py} (15).

\end{document}
 into a larger document.
%
%  Every table and section is \label'd for cross-referencing, using the
%  prefixes tab: and sec:. Numbers here are measured, not estimated;
%  where a number is an assumption it is marked as such.
% =====================================================================
\documentclass[11pt]{article}

\usepackage[margin=1in]{geometry}
\usepackage{amsmath,amssymb}
\usepackage{booktabs}
\usepackage{longtable}
\usepackage{siunitx}
\usepackage{hyperref}
\usepackage{xcolor}

\newcommand{\Pc}{P_c}
\newcommand{\worked}[1]{\textcolor{teal}{\textbf{Worked:}} #1}
\newcommand{\failed}[1]{\textcolor{red}{\textbf{Did not work:}} #1}
\newcommand{\gotcha}[1]{\textcolor{orange}{\textbf{Gotcha:}} #1}

\title{Surrogate-Accelerated Conjunction Assessment\\
       \large Running project log and results reference}
\author{Anurag Paudel}
\date{Last updated: 8 September 2026}

\begin{document}
\maketitle

\begin{abstract}
Working log for a project that estimates satellite collision probability
$\Pc$ using a Gaussian-process surrogate trained at space-filling design
points, benchmarked against a brute-force Monte Carlo baseline. This
document records what was built, what was measured, what failed, and which
claims are verified versus assumed. It is a reference sheet for the
write-up, not the write-up itself.
\end{abstract}

\tableofcontents

% =====================================================================
\section{Summary of headline results}\label{sec:headline}

\begin{longtable}{p{0.42\textwidth} p{0.5\textwidth}}
\toprule
\textbf{Result} & \textbf{Value} \\
\midrule
\endhead
TLE uncertainty, calibrated from data &
  $\sigma_R = \SI{95}{\meter}$, $\sigma_T = \SI{953}{\meter}$,
  $\sigma_N = \SI{143}{\meter}$ \\
Monte Carlo cost, $\Pc \approx 4\times10^{-6}$ &
  $\sim$\num{22000000} draws for \SI{10}{\percent} relative error;
  $\sim2.2\times10^{9}$ for \SI{1}{\percent} \\
Monte Carlo at \num{100000} draws & \textbf{zero hits} -- no estimate at all \\
Dynamic range of $\Pc$ over the parameter box & \num{12.8} orders of magnitude \\
Fraction of design points brute-force MC can label at $10^8$ draws &
  \SI{28}{\percent} \\
Dimensional reduction & $6 \to 4$ parameters, exact to $\sim10^{-12}$ \\
Design points saved by the reduction & $8$--$14\times$ at equal fill distance \\
Parallel tempering vs.\ best-of-20 random LHD & $2.1$--$3.1\times$ better $\psi$ \\
\bottomrule
\end{longtable}

% =====================================================================
\section{Phase 0: environment}\label{sec:phase0}

Python \num{3.9.6} (macOS system interpreter), virtual environment in
\texttt{venv/}. Dependencies: \texttt{sgp4}, \texttt{numpy},
\texttt{pandas}, \texttt{matplotlib}, \texttt{requests}, later
\texttt{scipy} and \texttt{pybind11}.

\gotcha{Python 3.9 is end-of-life, so \texttt{pip} resolved older builds
(\texttt{numpy 2.0.2}, \texttt{matplotlib 3.9.4}). The system \texttt{ssl}
module is compiled against LibreSSL 2.8.3, so every network call emits a
\texttt{NotOpenSSLWarning}. Harmless but noisy; a rebuild on Python
3.11+ would remove both issues.}

% =====================================================================
\section{Phase 1: data ingestion and propagation}\label{sec:phase1}

\subsection{TLE format}

Two-line element sets are fixed-column records, 69 characters per line.
Fields are located by column position, not delimiter. The six orbital
elements live on line 2: inclination, RAAN, eccentricity, argument of
perigee, mean anomaly, and mean motion (which stands in for the
semi-major axis, recovered as $a = (\mu/n^2)^{1/3}$).

Four format traps, all of which the parser handles explicitly:
\begin{enumerate}
  \item \textbf{Implied decimal points.} Eccentricity \texttt{0006703}
        means \num{0.0006703}. The $B^*$ and second-derivative fields use
        assumed-decimal exponential notation: \texttt{30074-3} is
        $0.30074\times10^{-3}$.
  \item \textbf{Two-digit epoch year.} 57--99 map to 19xx, 00--56 to 20xx.
  \item \textbf{Elements are Brouwer mean elements}, defined by the SGP4
        theory itself. Propagating them with a two-body Kepler propagator
        produces kilometres of error immediately. The TLE and SGP4 are a
        matched pair.
  \item \textbf{Output frame is TEME}, not J2000 and not ECEF.
\end{enumerate}

\subsection{What worked}

\worked{The ISS validation passed every check. Propagating three days
forward from epoch at one-minute cadence gave a mean altitude of
\SI{419.4}{\kilo\meter} (range \num{409.3}--\SI{427.7}{\kilo\meter}),
speed \num{7.645}--\SI{7.666}{\kilo\meter\per\second}, and a period of
\SI{92.96}{\minute} from the elements. An \emph{independent} empirical
period, measured by counting geocentric-radius minima over the
propagation, gave \SI{92.91}{\minute} -- agreement to
\SI{0.05}{\minute}, confirming the propagator is not simply echoing its
input.}

\subsection{Problems encountered}

\failed{My hand-typed ISS test fixture failed its own checksum on both
lines. The checksum implementation was correct -- it validated against
live CelesTrak data on both lines -- so the fixture's two check digits
were recomputed. Worth recording as a methodological point: when a test
fails, verify the test before ``fixing'' the code.}

\gotcha{CelesTrak returns HTTP 200 with a plaintext message, not an error
status, when a query matches nothing or the client is throttled. The
fetcher content-checks the response rather than trusting the status code.}

\gotcha{Space-Track registration completed \emph{instantly}, contrary to
documentation describing manual review taking days. Worth noting because
project planning had treated it as the long-lead dependency.}

% =====================================================================
\section{Phase 2: Monte Carlo baseline}\label{sec:phase2}

\subsection{What the public data actually contains}

The \texttt{cdm\_public} feed has \textbf{16 fields and no covariance
matrices}. It is a screening summary, not a full CCSDS Conjunction Data
Message: no state vectors, no epochs, no RTN covariance. Available:
\texttt{CDM\_ID}, \texttt{TCA}, \texttt{PC}, \texttt{MIN\_RNG},
\texttt{EMERGENCY\_REPORTABLE}, object IDs, names, types, RCS size class,
and exclusion volumes.

Three consequences shaped everything downstream:
\begin{enumerate}
  \item Covariances must be \textbf{synthesized}. This is a stated
        modelling assumption, not a data product.
  \item \texttt{PC} is the 18th Space Defense Squadron's own published
        probability. It is a reference point, \emph{not} a validation
        target, because it was computed from the covariance we do not
        have.
  \item The feed is pre-filtered to high-risk events: every sampled row
        had \texttt{EMERGENCY\_REPORTABLE = Y} and $\texttt{PC} > 10^{-4}$.
        There are \textbf{no negative examples}, which constrains the
        Phase 4 classifier.
\end{enumerate}

\gotcha{\texttt{EXCL\_VOL} is a kilometre-scale \emph{screening} volume
keyed to object class (debris 1, rocket body 3, payload 5), \textbf{not} a
hard-body radius. Since $\Pc \propto \text{HBR}^2$, using it as a hard-body
radius would have inflated every probability by roughly six orders of
magnitude. Values are fully determined by object type, which is what gave
it away.}

\subsection{Rebuilding real conjunctions from TLEs}

Both objects of each event were propagated to TCA with SGP4 and the
separation compared against the reported \texttt{MIN\_RNG}
(Table~\ref{tab:rebuild}).

\begin{table}[htbp]
\centering
\caption{Rebuilt versus reported miss distance, 37 deduplicated events.}
\label{tab:rebuild}
\begin{tabular}{lr}
\toprule
Quantity & Value (\si{\kilo\meter}) \\
\midrule
Reported miss distance (\texttt{MIN\_RNG}) range & 0.008--4.510 \\
Median $|$rebuilt $-$ reported$|$ & 0.936 \\
Mean $|$rebuilt $-$ reported$|$ & 2.219 \\
Range of $|$rebuilt $-$ reported$|$ & 0.046--20.798 \\
\bottomrule
\end{tabular}
\end{table}

This gap is the central motivating result: \textbf{reported miss distances
are hundreds of metres, while TLE-based rebuilds disagree by kilometres}.
Collision assessment cannot be done by propagating TLEs and measuring a
distance, because the error is an order of magnitude larger than the
quantity being measured. That is why the problem is probabilistic.

\subsection{Two observations worth keeping}

\textbf{Every rebuild error was positive.} The rebuilt miss is always
\emph{larger} than reported, never smaller. This is a genuine statistical
effect rather than a bias in the code: when the true miss is near zero,
adding random error to either object almost always increases the
separation, because distance is a norm. The Monte Carlo reproduces the
same asymmetry.

\textbf{Events are filed multiple times.} Each conjunction appears twice
(once with each object as primary) and is re-filed as the estimate is
refined, each revision at a slightly different TCA. Deduplication
therefore groups by unordered object pair and merges revisions within a
\SI{15}{\minute} window. An exact-TCA key silently missed the revisions.

\subsection{Covariance model (Option B)}

Diagonal in each object's RTN frame,
\begin{equation}
  C_{\text{RTN}} = \operatorname{diag}\!\left(\sigma_R^2,\;
    (k_T\sigma_R)^2,\; (k_N\sigma_R)^2\right),
  \label{eq:covb}
\end{equation}
with anisotropy ratios \emph{fixed} at $k_T = 10$, $k_N = 1.5$ and only
the scale $\sigma_R$ fitted. Twelve free parameters cannot be identified
from one scalar observable per event; one can. In-track error dominates
because a small period error integrates into a growing along-track lag.

Each object's covariance is rotated into the inertial frame \emph{before}
summing. The two objects occupy different orbital planes, so adding them
componentwise in RTN would be wrong.

\begin{table}[htbp]
\centering
\caption{Calibrated covariance scale, 37 events, ratios held fixed.}
\label{tab:calib}
\begin{tabular}{lS[table-format=4.0]}
\toprule
Component & {Value (\si{\meter})} \\
\midrule
$\sigma_R$ (radial)      & 95 \\
$\sigma_T$ (in-track)    & 953 \\
$\sigma_N$ (cross-track) & 143 \\
\bottomrule
\end{tabular}
\end{table}

\worked{These magnitudes are consistent with published TLE accuracy
($\sim$\SI{1}{\kilo\meter} in-track at a few days) and were derived from
the project's own measurements, not assumed.}

\failed{Per-event correlation between predicted and observed error is
$-0.53$ -- \emph{negative}. The model captures how large TLE error is in
aggregate but gets the per-event ranking backwards. The honest claim is
that it reproduces the scale of uncertainty, not which conjunctions are
worst. Cause: a single scalar observable per event is dominated by the
reported miss magnitude, and the model's per-event variation keys off the
same quantity in the opposite direction.}

\failed{Residual-versus-TLE-age correlation is only $-0.16$, so the data
does \emph{not} support an explicit time-growth term. Staying with the
fixed-scale model is an empirical decision, not a simplification.}

\subsection{Monte Carlo and its cost}

The short-term encounter model applies: at
\SIrange{10}{15}{\kilo\meter\per\second} relative speed the encounter
lasts milliseconds, relative motion is effectively rectilinear, and the
three-dimensional problem collapses onto the two-dimensional encounter
plane perpendicular to the relative velocity. $\Pc$ is then the integral
of the projected Gaussian over a disk of the combined hard-body radius.

Three independent estimators were implemented: batched Monte Carlo, polar
quadrature, and a small-disk closed form valid when
$\text{HBR} \ll \sigma$.

\worked{Monte Carlo and quadrature agree within Monte Carlo error bars on
every real event, maximum $|z| = 1.4$. Two methods sharing no code path,
so the encounter geometry and the sampler are cross-validated.}

\begin{table}[htbp]
\centering
\caption{Brute-force Monte Carlo convergence at $\Pc = 4.5\times10^{-6}$.
         This is the cost the surrogate replaces.}
\label{tab:mccost}
\begin{tabular}{rlS[table-format=3.1]}
\toprule
{Draws} & {$\Pc$ estimate} & {Relative error (\si{\percent})} \\
\midrule
\num{10000}     & \num{0.0}          & 100.0 \\
\num{100000}    & \num{0.0}          & 100.0 \\
\num{1000000}   & \num{6.0e-6}       & 33.1 \\
\num{10000000}  & \num{2.7e-6}       & 40.1 \\
\num{50000000}  & \num{4.56e-6}      & 1.2 \\
\bottomrule
\end{tabular}
\end{table}

At \num{100000} draws the estimator returns \textbf{zero hits} -- not a
poor answer, but no answer. Reaching \SI{10}{\percent} relative error
requires $\sim$\num{22000000} draws and \SI{1}{\percent} requires
$\sim2.2\times10^{9}$, per conjunction, against a catalogue of thousands
of screened pairs.

\subsection{Known caveats}\label{sec:caveats}

\begin{enumerate}
  \item Computed $\Pc$ runs $10$--$100\times$ \emph{below} the published
        \texttt{PC}. The dominant driver is almost certainly the hard-body
        radius: RCS-derived values of \SIrange{1}{4}{\meter} are likely
        below operational values, and $\Pc \propto \text{HBR}^2$, so a
        $10\times$ error in HBR is $100\times$ in $\Pc$. Agreement with
        \texttt{PC} must \emph{not} be presented as validation.
  \item Per-event calibration correlation is negative (above).
  \item One event showed a \SI{20.8}{\kilo\meter} rebuild error, far
        outside plausible TLE error. Likely a manoeuvre or stale elements;
        should be identified before it contaminates Phase 4.
  \item Slow encounters (relative velocity below
        \SI{1}{\kilo\meter\per\second}) violate the short-term encounter
        assumption and are currently filtered out.
\end{enumerate}

% =====================================================================
\section{Phase 3: space-filling designs and the surrogate}\label{sec:phase3}

\subsection{A negative result that redirected the phase}\label{sec:negative}

The intuitive reading of ``use a space-filling design instead of Monte
Carlo'' does \textbf{not} work. $\Pc$ is a rare-event indicator integral
over a microscopic disk. A 256-point design returns zero hits just as
\num{100000} random draws do, and in fact more reliably, because a
space-filling design deliberately spreads points \emph{away} from any
small region. Quasi-Monte Carlo does not rescue this either: its
advantage requires a smooth integrand, and an indicator function over a
tiny region is the opposite.

The construction that does work is a surrogate over the
\textbf{encounter-parameter space}, trained on expensive evaluations at
design points and amortised across the catalogue. This is recorded because
it is the kind of dead end worth reporting in a write-up rather than
quietly discarding.

\subsection{Dimensional reduction, verified rather than assumed}

$\Pc$ depends on six raw parameters: the miss vector $\mu$ (2), the
projected covariance $C$ (3), and the hard-body radius $R$ (1). Only
\textbf{four} are independent, because the defining integral
\begin{equation}
  \Pc = \int_{\|x\| \le R} \mathcal{N}(x;\, \mu,\, C)\, \mathrm{d}x
  \label{eq:pcdef}
\end{equation}
has a rotation-symmetric domain and is covariant under uniform scaling of
all lengths. Rotating into the eigenbasis of $C$ removes one parameter;
dividing every length by $R$ removes another, leaving
$(\mu_1/R,\, \mu_2/R,\, \sigma_1/R,\, \sigma_2/R)$.

This is an exact identity, and it was verified numerically rather than
argued from algebra (Table~\ref{tab:invariance}).

\begin{table}[htbp]
\centering
\caption{Numerical verification of the invariances underlying the
         $6 \to 4$ reduction. Errors are quadrature noise, not model
         error.}
\label{tab:invariance}
\begin{tabular}{lrS[table-format=1.1e-2]}
\toprule
Claim & {Cases} & {Max.\ relative error} \\
\midrule
$\Pc$ invariant under rotation of $\mu$ and $C$ & 300 & 8.3e-13 \\
$\Pc$ invariant under uniform length scaling    & 300 & 9.8e-13 \\
4-D reduction reproduces $\Pc$                  & 200 & 8.4e-13 \\
$\Pc$ even in each $\mu$ component (eigenbasis) & 200 & 6.2e-16 \\
$6 \to 4 \to \Pc$ round trip                    & 200 & 5.5e-13 \\
\bottomrule
\end{tabular}
\end{table}

\worked{The reduction is exact to machine precision, so it costs no
accuracy and saves substantial cost. Measured fill distance at $k=4$,
$n=64$ is \num{0.397}; matching that in six dimensions requires
$\sim$\numrange{500}{900} points, i.e.\ $8$--$14\times$ more expensive
evaluations for identical answers.}

The identity holds \emph{given the model}: circular hard body, Gaussian
uncertainty, short-term encounter. Slow encounters fall outside it. The
six-dimensional space is retained as an ablation so the saving can be
demonstrated empirically rather than asserted.

\subsection{The 4-D parameter space}

\begin{table}[htbp]
\centering
\caption{Primary parameter space. Ranges from 40 rebuilt conjunctions,
         widened. Anisotropy is parameterised as a ratio so
         $\sigma_2 \le \sigma_1$ holds by construction.}
\label{tab:paramspace}
\begin{tabular}{llll}
\toprule
Index & Parameter & Meaning & Range \\
\midrule
0 & $\log_{10}(\sigma_1/R)$        & uncertainty vs.\ object size & 32 -- 3162 \\
1 & $\log_{10}(\sigma_2/\sigma_1)$ & anisotropy                   & 0.032 -- 1 \\
2 & $\mu_1/\sigma_1$               & miss along major axis        & 0 -- 6 \\
3 & $\mu_2/\sigma_2$               & miss along minor axis        & 0 -- 6 \\
\bottomrule
\end{tabular}
\end{table}

\subsection{Design generation}

The RUSIS parallel-tempering MaxPro code was reused directly. The
pybind11 port compiled without modification and its designs come out
already scaled to $[0,1]$ as $(\ell - 0.5)/n$.

\worked{An independent \texttt{numpy} reimplementation of the MaxPro
criterion $\psi$ agrees with the C\texttt{++} extension to six decimal
places, cross-checking the original research code.}

\begin{table}[htbp]
\centering
\caption{Design quality at $k=6$. Lower $\psi$ and $\phi_p$ are better;
         larger minimum distance is better. Random LHD is the best of 20
         draws.}
\label{tab:designquality}
\begin{tabular}{rlSSS}
\toprule
{$n$} & Design & {$\psi$ (MaxPro)} & {$\phi_p$} & {Min.\ distance} \\
\midrule
64  & MaxPro (PT) & 37.7297 & 1.7649 & 0.3694 \\
64  & Random LHD  & 81.2452 & 2.3531 & 0.2684 \\
64  & Uniform     & 508.8451 & 2.8862 & 0.2118 \\
\midrule
256 & MaxPro (PT) & 79.6778 & 2.4753 & 0.2162 \\
256 & Random LHD  & 245.0499 & 3.4067 & 0.1533 \\
256 & Uniform     & 415.3693 & 3.3181 & 0.1615 \\
\bottomrule
\end{tabular}
\end{table}

Parallel tempering beats the best of 20 random Latin hypercubes by
$2.1$--$3.1\times$ on $\psi$, at \SIrange{0.2}{1.4}{\second} per design --
a cost paid once, since a design is a static artefact.

\subsection{Why training labels come from quadrature, not Monte Carlo}
\label{sec:labels}

$\Pc$ spans \textbf{12.8 orders of magnitude} across the parameter box
($2.1\times10^{-17}$ to $1.3\times10^{-4}$ on the $n=64$ design). Brute
force cannot label it:

\begin{table}[htbp]
\centering
\caption{Fraction of the 64 design points at which brute-force Monte Carlo
         yields a usable estimate (more than 10 expected hits).}
\label{tab:mcunusable}
\begin{tabular}{rr}
\toprule
{Draws per point} & {Points labelled (of 64)} \\
\midrule
$10^6$ & 6 \phantom{0}(\SI{9}{\percent}) \\
$10^7$ & 13 (\SI{20}{\percent}) \\
$10^8$ & 18 (\SI{28}{\percent}) \\
\bottomrule
\end{tabular}
\end{table}

Nineteen of 64 points lie below $\Pc = 10^{-12}$, which no Monte Carlo
will ever reach. Labels therefore come from the polar quadrature, which is
exact here and agrees with Monte Carlo wherever Monte Carlo works at all
(Section~\ref{sec:phase2}).

This strengthens rather than weakens the motivation: brute-force Monte
Carlo is not merely \emph{expensive} for rare events, it is
\textbf{unusable} across most of the parameter space. The surrogate
produces estimates where brute force produces nothing. The quadrature
serves as exact ground truth in this Gaussian testbed; in the general case
(non-Gaussian uncertainty, long-duration encounters, non-circular bodies)
no closed form exists and the expensive sampler is the only option.

\subsection{Surrogate}

Gaussian process with an anisotropic (ARD) squared-exponential kernel,
hyperparameters by maximising the log marginal likelihood. Targets are
$\log_{10}\Pc$ rather than $\Pc$: with 13 orders of magnitude of dynamic
range, absolute error is meaningless and a GP on the raw value would be
dominated by a handful of large-$\Pc$ points. The GP is plain
\texttt{numpy} (one Cholesky factorisation, two triangular solves); only
the optimiser comes from \texttt{scipy}.

% --- RESULTS TABLE INSERTED BELOW WHEN THE BENCHMARK RUN COMPLETES ---
% Included by project_log.tex. Regenerate the numbers with
%   ./venv/bin/python src/run_surrogate.py

\begin{table}[htbp]
\centering
\caption{Surrogate accuracy against a 3000-point held-out test set.
         RMSE is in $\log_{10}\Pc$, i.e.\ orders of magnitude. Every design
         type gets five replicates: parallel tempering is seeded from
         \texttt{random\_device}, so repeated runs give genuinely different
         designs. Spread is the standard deviation across replicates.}
\label{tab:surrogate}
\begin{tabular}{rlSSS}
\toprule
{$n$} & Design & {RMSE} & {Spread} & {Fit time (\si{\second})} \\
\midrule
32  & MaxPro (PT) & 0.034 & 0.001 & 0.15 \\
32  & Random LHD  & 0.046 & 0.024 & 0.18 \\
32  & Uniform     & 0.043 & 0.015 & 0.19 \\
\midrule
64  & MaxPro (PT) & 0.029 & 0.001 & 0.39 \\
64  & Random LHD  & 0.033 & 0.009 & 0.37 \\
64  & Uniform     & 0.045 & 0.013 & 0.28 \\
\midrule
128 & MaxPro (PT) & 0.018 & 0.003 & 0.89 \\
128 & Random LHD  & 0.022 & 0.004 & 1.06 \\
128 & Uniform     & 0.033 & 0.014 & 1.20 \\
\midrule
256 & MaxPro (PT) & 0.006 & 0.003 & 4.21 \\
256 & Random LHD  & 0.006 & 0.001 & 4.73 \\
256 & Uniform     & 0.005 & 0.001 & 5.32 \\
\bottomrule
\end{tabular}
\end{table}

\subsubsection*{Reading the benchmark honestly}

\worked{The surrogate is accurate in absolute terms. At $n=64$ the RMSE is
\num{0.029} orders of magnitude -- roughly \SI{7}{\percent} in $\Pc$ --
from 64 training evaluations, against a quantity spanning 13 orders of
magnitude. A single prediction costs \SI{8.3}{\micro\second} against
$\sim$\num{22000000} Monte Carlo draws for \SI{10}{\percent} accuracy at
one representative event.}

\worked{\textbf{The clearest advantage of parallel tempering is variance,
not mean accuracy.} At $n=32$ MaxPro's RMSE spread across replicates is
$\pm\num{0.001}$ against $\pm\num{0.024}$ for random Latin hypercubes --
more than twenty times more consistent. A random design is a lottery: some
draws are as good as an optimised one, others are much worse. With a
one-shot expensive evaluation budget, reproducibility is worth as much as
the mean improvement, and this is the practitioner-facing argument for the
method.}

\failed{\textbf{The $2$--$3\times$ advantage in the $\psi$ criterion does
not translate into a $2$--$3\times$ accuracy advantage.} Mean RMSE
improves by only $1.1$--$1.8\times$ over random LHD and uniform sampling
at $n \le 128$, and at $n=256$ the advantage disappears entirely
($0.99\times$ vs.\ random LHD, $0.74\times$ vs.\ uniform). Better
space-filling by the criterion is not the same as better surrogate
accuracy, and this should be stated plainly rather than glossed.}

\failed{At $n=256$ every design type converges to RMSE $\approx
\num{0.006}$. The target function $\log_{10}\Pc$ is smooth and
low-complexity, so once the box is adequately covered the design choice
stops mattering. \textbf{The interesting regime is small $n$}, which is
also the regime that matters when each evaluation is expensive.}

\gotcha{An earlier version of this benchmark compared one parallel-tempering
design against \emph{one} random draw, and reported MaxPro losing at
$n=128$ ($0.94\times$). That was sampling noise in the baseline, not a real
effect. Comparing a deterministic method against a single random
realisation is not a valid comparison; all figures above use five
replicates per design type.}

\subsubsection*{Points needed to match MaxPro at $n=64$}

Taking MaxPro at $n=64$ (RMSE \num{0.029}) as the target, random Latin
hypercubes first reach it at $n=128$ and uniform sampling at $n=256$ --
nominally $2\times$ and $4\times$ the evaluation budget. Given the
replicate spreads in Table~\ref{tab:surrogate} and the coarse grid of
sizes tested, these factors should be read as indicative rather than
precise; a finer sweep of $n$ would be needed to state them with
confidence.


% =====================================================================
\section{Reproducing everything}\label{sec:repro}

\begin{longtable}{p{0.42\textwidth} p{0.5\textwidth}}
\toprule
\textbf{Command} & \textbf{Produces} \\
\midrule
\endhead
\texttt{python src/validate\_iss.py}    & ISS propagation sanity checks \\
\texttt{python src/probe\_cdm.py}       & \texttt{cdm\_public} field inventory \\
\texttt{python src/rebuild\_events.py}  & Table~\ref{tab:rebuild} \\
\texttt{python src/calibrate.py}        & Table~\ref{tab:calib} \\
\texttt{python src/run\_montecarlo.py}  & Table~\ref{tab:mccost} \\
\texttt{python src/run\_surrogate.py}   & Table~\ref{tab:surrogate} \\
\texttt{python tests/test\_paramspace.py} & Table~\ref{tab:invariance} \\
\bottomrule
\end{longtable}

Test suites: \texttt{tests/test\_tle.py} (31),
\texttt{tests/test\_covariance.py} (22),
\texttt{tests/test\_montecarlo.py} (18),
\texttt{tests/test\_paramspace.py} (15).

\end{document}
 into a larger document.
%
%  Every table and section is \label'd for cross-referencing, using the
%  prefixes tab: and sec:. Numbers here are measured, not estimated;
%  where a number is an assumption it is marked as such.
% =====================================================================
\documentclass[11pt]{article}

\usepackage[margin=1in]{geometry}
\usepackage{amsmath,amssymb}
\usepackage{booktabs}
\usepackage{longtable}
\usepackage{siunitx}
\usepackage{hyperref}
\usepackage{xcolor}

\newcommand{\Pc}{P_c}
\newcommand{\worked}[1]{\textcolor{teal}{\textbf{Worked:}} #1}
\newcommand{\failed}[1]{\textcolor{red}{\textbf{Did not work:}} #1}
\newcommand{\gotcha}[1]{\textcolor{orange}{\textbf{Gotcha:}} #1}

\title{Surrogate-Accelerated Conjunction Assessment\\
       \large Running project log and results reference}
\author{Anurag Paudel}
\date{Last updated: 8 September 2026}

\begin{document}
\maketitle

\begin{abstract}
Working log for a project that estimates satellite collision probability
$\Pc$ using a Gaussian-process surrogate trained at space-filling design
points, benchmarked against a brute-force Monte Carlo baseline. This
document records what was built, what was measured, what failed, and which
claims are verified versus assumed. It is a reference sheet for the
write-up, not the write-up itself.
\end{abstract}

\tableofcontents

% =====================================================================
\section{Summary of headline results}\label{sec:headline}

\begin{longtable}{p{0.42\textwidth} p{0.5\textwidth}}
\toprule
\textbf{Result} & \textbf{Value} \\
\midrule
\endhead
TLE uncertainty, calibrated from data &
  $\sigma_R = \SI{95}{\meter}$, $\sigma_T = \SI{953}{\meter}$,
  $\sigma_N = \SI{143}{\meter}$ \\
Monte Carlo cost, $\Pc \approx 4\times10^{-6}$ &
  $\sim$\num{22000000} draws for \SI{10}{\percent} relative error;
  $\sim2.2\times10^{9}$ for \SI{1}{\percent} \\
Monte Carlo at \num{100000} draws & \textbf{zero hits} -- no estimate at all \\
Dynamic range of $\Pc$ over the parameter box & \num{12.8} orders of magnitude \\
Fraction of design points brute-force MC can label at $10^8$ draws &
  \SI{28}{\percent} \\
Dimensional reduction & $6 \to 4$ parameters, exact to $\sim10^{-12}$ \\
Design points saved by the reduction & $8$--$14\times$ at equal fill distance \\
Parallel tempering vs.\ best-of-20 random LHD & $2.1$--$3.1\times$ better $\psi$ \\
\bottomrule
\end{longtable}

% =====================================================================
\section{Phase 0: environment}\label{sec:phase0}

Python \num{3.9.6} (macOS system interpreter), virtual environment in
\texttt{venv/}. Dependencies: \texttt{sgp4}, \texttt{numpy},
\texttt{pandas}, \texttt{matplotlib}, \texttt{requests}, later
\texttt{scipy} and \texttt{pybind11}.

\gotcha{Python 3.9 is end-of-life, so \texttt{pip} resolved older builds
(\texttt{numpy 2.0.2}, \texttt{matplotlib 3.9.4}). The system \texttt{ssl}
module is compiled against LibreSSL 2.8.3, so every network call emits a
\texttt{NotOpenSSLWarning}. Harmless but noisy; a rebuild on Python
3.11+ would remove both issues.}

% =====================================================================
\section{Phase 1: data ingestion and propagation}\label{sec:phase1}

\subsection{TLE format}

Two-line element sets are fixed-column records, 69 characters per line.
Fields are located by column position, not delimiter. The six orbital
elements live on line 2: inclination, RAAN, eccentricity, argument of
perigee, mean anomaly, and mean motion (which stands in for the
semi-major axis, recovered as $a = (\mu/n^2)^{1/3}$).

Four format traps, all of which the parser handles explicitly:
\begin{enumerate}
  \item \textbf{Implied decimal points.} Eccentricity \texttt{0006703}
        means \num{0.0006703}. The $B^*$ and second-derivative fields use
        assumed-decimal exponential notation: \texttt{30074-3} is
        $0.30074\times10^{-3}$.
  \item \textbf{Two-digit epoch year.} 57--99 map to 19xx, 00--56 to 20xx.
  \item \textbf{Elements are Brouwer mean elements}, defined by the SGP4
        theory itself. Propagating them with a two-body Kepler propagator
        produces kilometres of error immediately. The TLE and SGP4 are a
        matched pair.
  \item \textbf{Output frame is TEME}, not J2000 and not ECEF.
\end{enumerate}

\subsection{What worked}

\worked{The ISS validation passed every check. Propagating three days
forward from epoch at one-minute cadence gave a mean altitude of
\SI{419.4}{\kilo\meter} (range \num{409.3}--\SI{427.7}{\kilo\meter}),
speed \num{7.645}--\SI{7.666}{\kilo\meter\per\second}, and a period of
\SI{92.96}{\minute} from the elements. An \emph{independent} empirical
period, measured by counting geocentric-radius minima over the
propagation, gave \SI{92.91}{\minute} -- agreement to
\SI{0.05}{\minute}, confirming the propagator is not simply echoing its
input.}

\subsection{Problems encountered}

\failed{My hand-typed ISS test fixture failed its own checksum on both
lines. The checksum implementation was correct -- it validated against
live CelesTrak data on both lines -- so the fixture's two check digits
were recomputed. Worth recording as a methodological point: when a test
fails, verify the test before ``fixing'' the code.}

\gotcha{CelesTrak returns HTTP 200 with a plaintext message, not an error
status, when a query matches nothing or the client is throttled. The
fetcher content-checks the response rather than trusting the status code.}

\gotcha{Space-Track registration completed \emph{instantly}, contrary to
documentation describing manual review taking days. Worth noting because
project planning had treated it as the long-lead dependency.}

% =====================================================================
\section{Phase 2: Monte Carlo baseline}\label{sec:phase2}

\subsection{What the public data actually contains}

The \texttt{cdm\_public} feed has \textbf{16 fields and no covariance
matrices}. It is a screening summary, not a full CCSDS Conjunction Data
Message: no state vectors, no epochs, no RTN covariance. Available:
\texttt{CDM\_ID}, \texttt{TCA}, \texttt{PC}, \texttt{MIN\_RNG},
\texttt{EMERGENCY\_REPORTABLE}, object IDs, names, types, RCS size class,
and exclusion volumes.

Three consequences shaped everything downstream:
\begin{enumerate}
  \item Covariances must be \textbf{synthesized}. This is a stated
        modelling assumption, not a data product.
  \item \texttt{PC} is the 18th Space Defense Squadron's own published
        probability. It is a reference point, \emph{not} a validation
        target, because it was computed from the covariance we do not
        have.
  \item The feed is pre-filtered to high-risk events: every sampled row
        had \texttt{EMERGENCY\_REPORTABLE = Y} and $\texttt{PC} > 10^{-4}$.
        There are \textbf{no negative examples}, which constrains the
        Phase 4 classifier.
\end{enumerate}

\gotcha{\texttt{EXCL\_VOL} is a kilometre-scale \emph{screening} volume
keyed to object class (debris 1, rocket body 3, payload 5), \textbf{not} a
hard-body radius. Since $\Pc \propto \text{HBR}^2$, using it as a hard-body
radius would have inflated every probability by roughly six orders of
magnitude. Values are fully determined by object type, which is what gave
it away.}

\subsection{Rebuilding real conjunctions from TLEs}

Both objects of each event were propagated to TCA with SGP4 and the
separation compared against the reported \texttt{MIN\_RNG}
(Table~\ref{tab:rebuild}).

\begin{table}[htbp]
\centering
\caption{Rebuilt versus reported miss distance, 37 deduplicated events.}
\label{tab:rebuild}
\begin{tabular}{lr}
\toprule
Quantity & Value (\si{\kilo\meter}) \\
\midrule
Reported miss distance (\texttt{MIN\_RNG}) range & 0.008--4.510 \\
Median $|$rebuilt $-$ reported$|$ & 0.936 \\
Mean $|$rebuilt $-$ reported$|$ & 2.219 \\
Range of $|$rebuilt $-$ reported$|$ & 0.046--20.798 \\
\bottomrule
\end{tabular}
\end{table}

This gap is the central motivating result: \textbf{reported miss distances
are hundreds of metres, while TLE-based rebuilds disagree by kilometres}.
Collision assessment cannot be done by propagating TLEs and measuring a
distance, because the error is an order of magnitude larger than the
quantity being measured. That is why the problem is probabilistic.

\subsection{Two observations worth keeping}

\textbf{Every rebuild error was positive.} The rebuilt miss is always
\emph{larger} than reported, never smaller. This is a genuine statistical
effect rather than a bias in the code: when the true miss is near zero,
adding random error to either object almost always increases the
separation, because distance is a norm. The Monte Carlo reproduces the
same asymmetry.

\textbf{Events are filed multiple times.} Each conjunction appears twice
(once with each object as primary) and is re-filed as the estimate is
refined, each revision at a slightly different TCA. Deduplication
therefore groups by unordered object pair and merges revisions within a
\SI{15}{\minute} window. An exact-TCA key silently missed the revisions.

\subsection{Covariance model (Option B)}

Diagonal in each object's RTN frame,
\begin{equation}
  C_{\text{RTN}} = \operatorname{diag}\!\left(\sigma_R^2,\;
    (k_T\sigma_R)^2,\; (k_N\sigma_R)^2\right),
  \label{eq:covb}
\end{equation}
with anisotropy ratios \emph{fixed} at $k_T = 10$, $k_N = 1.5$ and only
the scale $\sigma_R$ fitted. Twelve free parameters cannot be identified
from one scalar observable per event; one can. In-track error dominates
because a small period error integrates into a growing along-track lag.

Each object's covariance is rotated into the inertial frame \emph{before}
summing. The two objects occupy different orbital planes, so adding them
componentwise in RTN would be wrong.

\begin{table}[htbp]
\centering
\caption{Calibrated covariance scale, 37 events, ratios held fixed.}
\label{tab:calib}
\begin{tabular}{lS[table-format=4.0]}
\toprule
Component & {Value (\si{\meter})} \\
\midrule
$\sigma_R$ (radial)      & 95 \\
$\sigma_T$ (in-track)    & 953 \\
$\sigma_N$ (cross-track) & 143 \\
\bottomrule
\end{tabular}
\end{table}

\worked{These magnitudes are consistent with published TLE accuracy
($\sim$\SI{1}{\kilo\meter} in-track at a few days) and were derived from
the project's own measurements, not assumed.}

\failed{Per-event correlation between predicted and observed error is
$-0.53$ -- \emph{negative}. The model captures how large TLE error is in
aggregate but gets the per-event ranking backwards. The honest claim is
that it reproduces the scale of uncertainty, not which conjunctions are
worst. Cause: a single scalar observable per event is dominated by the
reported miss magnitude, and the model's per-event variation keys off the
same quantity in the opposite direction.}

\failed{Residual-versus-TLE-age correlation is only $-0.16$, so the data
does \emph{not} support an explicit time-growth term. Staying with the
fixed-scale model is an empirical decision, not a simplification.}

\subsection{Monte Carlo and its cost}

The short-term encounter model applies: at
\SIrange{10}{15}{\kilo\meter\per\second} relative speed the encounter
lasts milliseconds, relative motion is effectively rectilinear, and the
three-dimensional problem collapses onto the two-dimensional encounter
plane perpendicular to the relative velocity. $\Pc$ is then the integral
of the projected Gaussian over a disk of the combined hard-body radius.

Three independent estimators were implemented: batched Monte Carlo, polar
quadrature, and a small-disk closed form valid when
$\text{HBR} \ll \sigma$.

\worked{Monte Carlo and quadrature agree within Monte Carlo error bars on
every real event, maximum $|z| = 1.4$. Two methods sharing no code path,
so the encounter geometry and the sampler are cross-validated.}

\begin{table}[htbp]
\centering
\caption{Brute-force Monte Carlo convergence at $\Pc = 4.5\times10^{-6}$.
         This is the cost the surrogate replaces.}
\label{tab:mccost}
\begin{tabular}{rlS[table-format=3.1]}
\toprule
{Draws} & {$\Pc$ estimate} & {Relative error (\si{\percent})} \\
\midrule
\num{10000}     & \num{0.0}          & 100.0 \\
\num{100000}    & \num{0.0}          & 100.0 \\
\num{1000000}   & \num{6.0e-6}       & 33.1 \\
\num{10000000}  & \num{2.7e-6}       & 40.1 \\
\num{50000000}  & \num{4.56e-6}      & 1.2 \\
\bottomrule
\end{tabular}
\end{table}

At \num{100000} draws the estimator returns \textbf{zero hits} -- not a
poor answer, but no answer. Reaching \SI{10}{\percent} relative error
requires $\sim$\num{22000000} draws and \SI{1}{\percent} requires
$\sim2.2\times10^{9}$, per conjunction, against a catalogue of thousands
of screened pairs.

\subsection{Known caveats}\label{sec:caveats}

\begin{enumerate}
  \item Computed $\Pc$ runs $10$--$100\times$ \emph{below} the published
        \texttt{PC}. The dominant driver is almost certainly the hard-body
        radius: RCS-derived values of \SIrange{1}{4}{\meter} are likely
        below operational values, and $\Pc \propto \text{HBR}^2$, so a
        $10\times$ error in HBR is $100\times$ in $\Pc$. Agreement with
        \texttt{PC} must \emph{not} be presented as validation.
  \item Per-event calibration correlation is negative (above).
  \item One event showed a \SI{20.8}{\kilo\meter} rebuild error, far
        outside plausible TLE error. Likely a manoeuvre or stale elements;
        should be identified before it contaminates Phase 4.
  \item Slow encounters (relative velocity below
        \SI{1}{\kilo\meter\per\second}) violate the short-term encounter
        assumption and are currently filtered out.
\end{enumerate}

% =====================================================================
\section{Phase 3: space-filling designs and the surrogate}\label{sec:phase3}

\subsection{A negative result that redirected the phase}\label{sec:negative}

The intuitive reading of ``use a space-filling design instead of Monte
Carlo'' does \textbf{not} work. $\Pc$ is a rare-event indicator integral
over a microscopic disk. A 256-point design returns zero hits just as
\num{100000} random draws do, and in fact more reliably, because a
space-filling design deliberately spreads points \emph{away} from any
small region. Quasi-Monte Carlo does not rescue this either: its
advantage requires a smooth integrand, and an indicator function over a
tiny region is the opposite.

The construction that does work is a surrogate over the
\textbf{encounter-parameter space}, trained on expensive evaluations at
design points and amortised across the catalogue. This is recorded because
it is the kind of dead end worth reporting in a write-up rather than
quietly discarding.

\subsection{Dimensional reduction, verified rather than assumed}

$\Pc$ depends on six raw parameters: the miss vector $\mu$ (2), the
projected covariance $C$ (3), and the hard-body radius $R$ (1). Only
\textbf{four} are independent, because the defining integral
\begin{equation}
  \Pc = \int_{\|x\| \le R} \mathcal{N}(x;\, \mu,\, C)\, \mathrm{d}x
  \label{eq:pcdef}
\end{equation}
has a rotation-symmetric domain and is covariant under uniform scaling of
all lengths. Rotating into the eigenbasis of $C$ removes one parameter;
dividing every length by $R$ removes another, leaving
$(\mu_1/R,\, \mu_2/R,\, \sigma_1/R,\, \sigma_2/R)$.

This is an exact identity, and it was verified numerically rather than
argued from algebra (Table~\ref{tab:invariance}).

\begin{table}[htbp]
\centering
\caption{Numerical verification of the invariances underlying the
         $6 \to 4$ reduction. Errors are quadrature noise, not model
         error.}
\label{tab:invariance}
\begin{tabular}{lrS[table-format=1.1e-2]}
\toprule
Claim & {Cases} & {Max.\ relative error} \\
\midrule
$\Pc$ invariant under rotation of $\mu$ and $C$ & 300 & 8.3e-13 \\
$\Pc$ invariant under uniform length scaling    & 300 & 9.8e-13 \\
4-D reduction reproduces $\Pc$                  & 200 & 8.4e-13 \\
$\Pc$ even in each $\mu$ component (eigenbasis) & 200 & 6.2e-16 \\
$6 \to 4 \to \Pc$ round trip                    & 200 & 5.5e-13 \\
\bottomrule
\end{tabular}
\end{table}

\worked{The reduction is exact to machine precision, so it costs no
accuracy and saves substantial cost. Measured fill distance at $k=4$,
$n=64$ is \num{0.397}; matching that in six dimensions requires
$\sim$\numrange{500}{900} points, i.e.\ $8$--$14\times$ more expensive
evaluations for identical answers.}

The identity holds \emph{given the model}: circular hard body, Gaussian
uncertainty, short-term encounter. Slow encounters fall outside it. The
six-dimensional space is retained as an ablation so the saving can be
demonstrated empirically rather than asserted.

\subsection{The 4-D parameter space}

\begin{table}[htbp]
\centering
\caption{Primary parameter space. Ranges from 40 rebuilt conjunctions,
         widened. Anisotropy is parameterised as a ratio so
         $\sigma_2 \le \sigma_1$ holds by construction.}
\label{tab:paramspace}
\begin{tabular}{llll}
\toprule
Index & Parameter & Meaning & Range \\
\midrule
0 & $\log_{10}(\sigma_1/R)$        & uncertainty vs.\ object size & 32 -- 3162 \\
1 & $\log_{10}(\sigma_2/\sigma_1)$ & anisotropy                   & 0.032 -- 1 \\
2 & $\mu_1/\sigma_1$               & miss along major axis        & 0 -- 6 \\
3 & $\mu_2/\sigma_2$               & miss along minor axis        & 0 -- 6 \\
\bottomrule
\end{tabular}
\end{table}

\subsection{Design generation}

The RUSIS parallel-tempering MaxPro code was reused directly. The
pybind11 port compiled without modification and its designs come out
already scaled to $[0,1]$ as $(\ell - 0.5)/n$.

\worked{An independent \texttt{numpy} reimplementation of the MaxPro
criterion $\psi$ agrees with the C\texttt{++} extension to six decimal
places, cross-checking the original research code.}

\begin{table}[htbp]
\centering
\caption{Design quality at $k=6$. Lower $\psi$ and $\phi_p$ are better;
         larger minimum distance is better. Random LHD is the best of 20
         draws.}
\label{tab:designquality}
\begin{tabular}{rlSSS}
\toprule
{$n$} & Design & {$\psi$ (MaxPro)} & {$\phi_p$} & {Min.\ distance} \\
\midrule
64  & MaxPro (PT) & 37.7297 & 1.7649 & 0.3694 \\
64  & Random LHD  & 81.2452 & 2.3531 & 0.2684 \\
64  & Uniform     & 508.8451 & 2.8862 & 0.2118 \\
\midrule
256 & MaxPro (PT) & 79.6778 & 2.4753 & 0.2162 \\
256 & Random LHD  & 245.0499 & 3.4067 & 0.1533 \\
256 & Uniform     & 415.3693 & 3.3181 & 0.1615 \\
\bottomrule
\end{tabular}
\end{table}

Parallel tempering beats the best of 20 random Latin hypercubes by
$2.1$--$3.1\times$ on $\psi$, at \SIrange{0.2}{1.4}{\second} per design --
a cost paid once, since a design is a static artefact.

\subsection{Why training labels come from quadrature, not Monte Carlo}
\label{sec:labels}

$\Pc$ spans \textbf{12.8 orders of magnitude} across the parameter box
($2.1\times10^{-17}$ to $1.3\times10^{-4}$ on the $n=64$ design). Brute
force cannot label it:

\begin{table}[htbp]
\centering
\caption{Fraction of the 64 design points at which brute-force Monte Carlo
         yields a usable estimate (more than 10 expected hits).}
\label{tab:mcunusable}
\begin{tabular}{rr}
\toprule
{Draws per point} & {Points labelled (of 64)} \\
\midrule
$10^6$ & 6 \phantom{0}(\SI{9}{\percent}) \\
$10^7$ & 13 (\SI{20}{\percent}) \\
$10^8$ & 18 (\SI{28}{\percent}) \\
\bottomrule
\end{tabular}
\end{table}

Nineteen of 64 points lie below $\Pc = 10^{-12}$, which no Monte Carlo
will ever reach. Labels therefore come from the polar quadrature, which is
exact here and agrees with Monte Carlo wherever Monte Carlo works at all
(Section~\ref{sec:phase2}).

This strengthens rather than weakens the motivation: brute-force Monte
Carlo is not merely \emph{expensive} for rare events, it is
\textbf{unusable} across most of the parameter space. The surrogate
produces estimates where brute force produces nothing. The quadrature
serves as exact ground truth in this Gaussian testbed; in the general case
(non-Gaussian uncertainty, long-duration encounters, non-circular bodies)
no closed form exists and the expensive sampler is the only option.

\subsection{Surrogate}

Gaussian process with an anisotropic (ARD) squared-exponential kernel,
hyperparameters by maximising the log marginal likelihood. Targets are
$\log_{10}\Pc$ rather than $\Pc$: with 13 orders of magnitude of dynamic
range, absolute error is meaningless and a GP on the raw value would be
dominated by a handful of large-$\Pc$ points. The GP is plain
\texttt{numpy} (one Cholesky factorisation, two triangular solves); only
the optimiser comes from \texttt{scipy}.

% --- RESULTS TABLE INSERTED BELOW WHEN THE BENCHMARK RUN COMPLETES ---
% Included by project_log.tex. Regenerate the numbers with
%   ./venv/bin/python src/run_surrogate.py

\begin{table}[htbp]
\centering
\caption{Surrogate accuracy against a 3000-point held-out test set.
         RMSE is in $\log_{10}\Pc$, i.e.\ orders of magnitude. Every design
         type gets five replicates: parallel tempering is seeded from
         \texttt{random\_device}, so repeated runs give genuinely different
         designs. Spread is the standard deviation across replicates.}
\label{tab:surrogate}
\begin{tabular}{rlSSS}
\toprule
{$n$} & Design & {RMSE} & {Spread} & {Fit time (\si{\second})} \\
\midrule
32  & MaxPro (PT) & 0.034 & 0.001 & 0.15 \\
32  & Random LHD  & 0.046 & 0.024 & 0.18 \\
32  & Uniform     & 0.043 & 0.015 & 0.19 \\
\midrule
64  & MaxPro (PT) & 0.029 & 0.001 & 0.39 \\
64  & Random LHD  & 0.033 & 0.009 & 0.37 \\
64  & Uniform     & 0.045 & 0.013 & 0.28 \\
\midrule
128 & MaxPro (PT) & 0.018 & 0.003 & 0.89 \\
128 & Random LHD  & 0.022 & 0.004 & 1.06 \\
128 & Uniform     & 0.033 & 0.014 & 1.20 \\
\midrule
256 & MaxPro (PT) & 0.006 & 0.003 & 4.21 \\
256 & Random LHD  & 0.006 & 0.001 & 4.73 \\
256 & Uniform     & 0.005 & 0.001 & 5.32 \\
\bottomrule
\end{tabular}
\end{table}

\subsubsection*{Reading the benchmark honestly}

\worked{The surrogate is accurate in absolute terms. At $n=64$ the RMSE is
\num{0.029} orders of magnitude -- roughly \SI{7}{\percent} in $\Pc$ --
from 64 training evaluations, against a quantity spanning 13 orders of
magnitude. A single prediction costs \SI{8.3}{\micro\second} against
$\sim$\num{22000000} Monte Carlo draws for \SI{10}{\percent} accuracy at
one representative event.}

\worked{\textbf{The clearest advantage of parallel tempering is variance,
not mean accuracy.} At $n=32$ MaxPro's RMSE spread across replicates is
$\pm\num{0.001}$ against $\pm\num{0.024}$ for random Latin hypercubes --
more than twenty times more consistent. A random design is a lottery: some
draws are as good as an optimised one, others are much worse. With a
one-shot expensive evaluation budget, reproducibility is worth as much as
the mean improvement, and this is the practitioner-facing argument for the
method.}

\failed{\textbf{The $2$--$3\times$ advantage in the $\psi$ criterion does
not translate into a $2$--$3\times$ accuracy advantage.} Mean RMSE
improves by only $1.1$--$1.8\times$ over random LHD and uniform sampling
at $n \le 128$, and at $n=256$ the advantage disappears entirely
($0.99\times$ vs.\ random LHD, $0.74\times$ vs.\ uniform). Better
space-filling by the criterion is not the same as better surrogate
accuracy, and this should be stated plainly rather than glossed.}

\failed{At $n=256$ every design type converges to RMSE $\approx
\num{0.006}$. The target function $\log_{10}\Pc$ is smooth and
low-complexity, so once the box is adequately covered the design choice
stops mattering. \textbf{The interesting regime is small $n$}, which is
also the regime that matters when each evaluation is expensive.}

\gotcha{An earlier version of this benchmark compared one parallel-tempering
design against \emph{one} random draw, and reported MaxPro losing at
$n=128$ ($0.94\times$). That was sampling noise in the baseline, not a real
effect. Comparing a deterministic method against a single random
realisation is not a valid comparison; all figures above use five
replicates per design type.}

\subsubsection*{Points needed to match MaxPro at $n=64$}

Taking MaxPro at $n=64$ (RMSE \num{0.029}) as the target, random Latin
hypercubes first reach it at $n=128$ and uniform sampling at $n=256$ --
nominally $2\times$ and $4\times$ the evaluation budget. Given the
replicate spreads in Table~\ref{tab:surrogate} and the coarse grid of
sizes tested, these factors should be read as indicative rather than
precise; a finer sweep of $n$ would be needed to state them with
confidence.


% =====================================================================
\section{Reproducing everything}\label{sec:repro}

\begin{longtable}{p{0.42\textwidth} p{0.5\textwidth}}
\toprule
\textbf{Command} & \textbf{Produces} \\
\midrule
\endhead
\texttt{python src/validate\_iss.py}    & ISS propagation sanity checks \\
\texttt{python src/probe\_cdm.py}       & \texttt{cdm\_public} field inventory \\
\texttt{python src/rebuild\_events.py}  & Table~\ref{tab:rebuild} \\
\texttt{python src/calibrate.py}        & Table~\ref{tab:calib} \\
\texttt{python src/run\_montecarlo.py}  & Table~\ref{tab:mccost} \\
\texttt{python src/run\_surrogate.py}   & Table~\ref{tab:surrogate} \\
\texttt{python tests/test\_paramspace.py} & Table~\ref{tab:invariance} \\
\bottomrule
\end{longtable}

Test suites: \texttt{tests/test\_tle.py} (31),
\texttt{tests/test\_covariance.py} (22),
\texttt{tests/test\_montecarlo.py} (18),
\texttt{tests/test\_paramspace.py} (15).

\end{document}
 into a larger document.
%
%  Every table and section is \label'd for cross-referencing, using the
%  prefixes tab: and sec:. Numbers here are measured, not estimated;
%  where a number is an assumption it is marked as such.
% =====================================================================
\documentclass[11pt]{article}

\usepackage[margin=1in]{geometry}
\usepackage{amsmath,amssymb}
\usepackage{booktabs}
\usepackage{longtable}
\usepackage{siunitx}
\usepackage{hyperref}
\usepackage{xcolor}

\newcommand{\Pc}{P_c}
\newcommand{\worked}[1]{\textcolor{teal}{\textbf{Worked:}} #1}
\newcommand{\failed}[1]{\textcolor{red}{\textbf{Did not work:}} #1}
\newcommand{\gotcha}[1]{\textcolor{orange}{\textbf{Gotcha:}} #1}

\title{Surrogate-Accelerated Conjunction Assessment\\
       \large Running project log and results reference}
\author{Anurag Paudel}
\date{Last updated: 8 September 2026}

\begin{document}
\maketitle

\begin{abstract}
Working log for a project that estimates satellite collision probability
$\Pc$ using a Gaussian-process surrogate trained at space-filling design
points, benchmarked against a brute-force Monte Carlo baseline. This
document records what was built, what was measured, what failed, and which
claims are verified versus assumed. It is a reference sheet for the
write-up, not the write-up itself.
\end{abstract}

\tableofcontents

% =====================================================================
\section{Summary of headline results}\label{sec:headline}

\begin{longtable}{p{0.42\textwidth} p{0.5\textwidth}}
\toprule
\textbf{Result} & \textbf{Value} \\
\midrule
\endhead
TLE uncertainty, calibrated from data &
  $\sigma_R = \SI{95}{\meter}$, $\sigma_T = \SI{953}{\meter}$,
  $\sigma_N = \SI{143}{\meter}$ \\
Monte Carlo cost, $\Pc \approx 4\times10^{-6}$ &
  $\sim$\num{22000000} draws for \SI{10}{\percent} relative error;
  $\sim2.2\times10^{9}$ for \SI{1}{\percent} \\
Monte Carlo at \num{100000} draws & \textbf{zero hits} -- no estimate at all \\
Dynamic range of $\Pc$ over the parameter box & \num{12.8} orders of magnitude \\
Fraction of design points brute-force MC can label at $10^8$ draws &
  \SI{28}{\percent} \\
Dimensional reduction & $6 \to 4$ parameters, exact to $\sim10^{-12}$ \\
Design points saved by the reduction & $8$--$14\times$ at equal fill distance \\
Parallel tempering vs.\ best-of-20 random LHD & $2.1$--$3.1\times$ better $\psi$ \\
\bottomrule
\end{longtable}

% =====================================================================
\section{Phase 0: environment}\label{sec:phase0}

Python \num{3.9.6} (macOS system interpreter), virtual environment in
\texttt{venv/}. Dependencies: \texttt{sgp4}, \texttt{numpy},
\texttt{pandas}, \texttt{matplotlib}, \texttt{requests}, later
\texttt{scipy} and \texttt{pybind11}.

\gotcha{Python 3.9 is end-of-life, so \texttt{pip} resolved older builds
(\texttt{numpy 2.0.2}, \texttt{matplotlib 3.9.4}). The system \texttt{ssl}
module is compiled against LibreSSL 2.8.3, so every network call emits a
\texttt{NotOpenSSLWarning}. Harmless but noisy; a rebuild on Python
3.11+ would remove both issues.}

% =====================================================================
\section{Phase 1: data ingestion and propagation}\label{sec:phase1}

\subsection{TLE format}

Two-line element sets are fixed-column records, 69 characters per line.
Fields are located by column position, not delimiter. The six orbital
elements live on line 2: inclination, RAAN, eccentricity, argument of
perigee, mean anomaly, and mean motion (which stands in for the
semi-major axis, recovered as $a = (\mu/n^2)^{1/3}$).

Four format traps, all of which the parser handles explicitly:
\begin{enumerate}
  \item \textbf{Implied decimal points.} Eccentricity \texttt{0006703}
        means \num{0.0006703}. The $B^*$ and second-derivative fields use
        assumed-decimal exponential notation: \texttt{30074-3} is
        $0.30074\times10^{-3}$.
  \item \textbf{Two-digit epoch year.} 57--99 map to 19xx, 00--56 to 20xx.
  \item \textbf{Elements are Brouwer mean elements}, defined by the SGP4
        theory itself. Propagating them with a two-body Kepler propagator
        produces kilometres of error immediately. The TLE and SGP4 are a
        matched pair.
  \item \textbf{Output frame is TEME}, not J2000 and not ECEF.
\end{enumerate}

\subsection{What worked}

\worked{The ISS validation passed every check. Propagating three days
forward from epoch at one-minute cadence gave a mean altitude of
\SI{419.4}{\kilo\meter} (range \num{409.3}--\SI{427.7}{\kilo\meter}),
speed \num{7.645}--\SI{7.666}{\kilo\meter\per\second}, and a period of
\SI{92.96}{\minute} from the elements. An \emph{independent} empirical
period, measured by counting geocentric-radius minima over the
propagation, gave \SI{92.91}{\minute} -- agreement to
\SI{0.05}{\minute}, confirming the propagator is not simply echoing its
input.}

\subsection{Problems encountered}

\failed{My hand-typed ISS test fixture failed its own checksum on both
lines. The checksum implementation was correct -- it validated against
live CelesTrak data on both lines -- so the fixture's two check digits
were recomputed. Worth recording as a methodological point: when a test
fails, verify the test before ``fixing'' the code.}

\gotcha{CelesTrak returns HTTP 200 with a plaintext message, not an error
status, when a query matches nothing or the client is throttled. The
fetcher content-checks the response rather than trusting the status code.}

\gotcha{Space-Track registration completed \emph{instantly}, contrary to
documentation describing manual review taking days. Worth noting because
project planning had treated it as the long-lead dependency.}

% =====================================================================
\section{Phase 2: Monte Carlo baseline}\label{sec:phase2}

\subsection{What the public data actually contains}

The \texttt{cdm\_public} feed has \textbf{16 fields and no covariance
matrices}. It is a screening summary, not a full CCSDS Conjunction Data
Message: no state vectors, no epochs, no RTN covariance. Available:
\texttt{CDM\_ID}, \texttt{TCA}, \texttt{PC}, \texttt{MIN\_RNG},
\texttt{EMERGENCY\_REPORTABLE}, object IDs, names, types, RCS size class,
and exclusion volumes.

Three consequences shaped everything downstream:
\begin{enumerate}
  \item Covariances must be \textbf{synthesized}. This is a stated
        modelling assumption, not a data product.
  \item \texttt{PC} is the 18th Space Defense Squadron's own published
        probability. It is a reference point, \emph{not} a validation
        target, because it was computed from the covariance we do not
        have.
  \item The feed is pre-filtered to high-risk events: every sampled row
        had \texttt{EMERGENCY\_REPORTABLE = Y} and $\texttt{PC} > 10^{-4}$.
        There are \textbf{no negative examples}, which constrains the
        Phase 4 classifier.
\end{enumerate}

\gotcha{\texttt{EXCL\_VOL} is a kilometre-scale \emph{screening} volume
keyed to object class (debris 1, rocket body 3, payload 5), \textbf{not} a
hard-body radius. Since $\Pc \propto \text{HBR}^2$, using it as a hard-body
radius would have inflated every probability by roughly six orders of
magnitude. Values are fully determined by object type, which is what gave
it away.}

\subsection{Rebuilding real conjunctions from TLEs}

Both objects of each event were propagated to TCA with SGP4 and the
separation compared against the reported \texttt{MIN\_RNG}
(Table~\ref{tab:rebuild}).

\begin{table}[htbp]
\centering
\caption{Rebuilt versus reported miss distance, 37 deduplicated events.}
\label{tab:rebuild}
\begin{tabular}{lr}
\toprule
Quantity & Value (\si{\kilo\meter}) \\
\midrule
Reported miss distance (\texttt{MIN\_RNG}) range & 0.008--4.510 \\
Median $|$rebuilt $-$ reported$|$ & 0.936 \\
Mean $|$rebuilt $-$ reported$|$ & 2.219 \\
Range of $|$rebuilt $-$ reported$|$ & 0.046--20.798 \\
\bottomrule
\end{tabular}
\end{table}

This gap is the central motivating result: \textbf{reported miss distances
are hundreds of metres, while TLE-based rebuilds disagree by kilometres}.
Collision assessment cannot be done by propagating TLEs and measuring a
distance, because the error is an order of magnitude larger than the
quantity being measured. That is why the problem is probabilistic.

\subsection{Two observations worth keeping}

\textbf{Every rebuild error was positive.} The rebuilt miss is always
\emph{larger} than reported, never smaller. This is a genuine statistical
effect rather than a bias in the code: when the true miss is near zero,
adding random error to either object almost always increases the
separation, because distance is a norm. The Monte Carlo reproduces the
same asymmetry.

\textbf{Events are filed multiple times.} Each conjunction appears twice
(once with each object as primary) and is re-filed as the estimate is
refined, each revision at a slightly different TCA. Deduplication
therefore groups by unordered object pair and merges revisions within a
\SI{15}{\minute} window. An exact-TCA key silently missed the revisions.

\subsection{Covariance model (Option B)}

Diagonal in each object's RTN frame,
\begin{equation}
  C_{\text{RTN}} = \operatorname{diag}\!\left(\sigma_R^2,\;
    (k_T\sigma_R)^2,\; (k_N\sigma_R)^2\right),
  \label{eq:covb}
\end{equation}
with anisotropy ratios \emph{fixed} at $k_T = 10$, $k_N = 1.5$ and only
the scale $\sigma_R$ fitted. Twelve free parameters cannot be identified
from one scalar observable per event; one can. In-track error dominates
because a small period error integrates into a growing along-track lag.

Each object's covariance is rotated into the inertial frame \emph{before}
summing. The two objects occupy different orbital planes, so adding them
componentwise in RTN would be wrong.

\begin{table}[htbp]
\centering
\caption{Calibrated covariance scale, 37 events, ratios held fixed.}
\label{tab:calib}
\begin{tabular}{lS[table-format=4.0]}
\toprule
Component & {Value (\si{\meter})} \\
\midrule
$\sigma_R$ (radial)      & 95 \\
$\sigma_T$ (in-track)    & 953 \\
$\sigma_N$ (cross-track) & 143 \\
\bottomrule
\end{tabular}
\end{table}

\worked{These magnitudes are consistent with published TLE accuracy
($\sim$\SI{1}{\kilo\meter} in-track at a few days) and were derived from
the project's own measurements, not assumed.}

\failed{Per-event correlation between predicted and observed error is
$-0.53$ -- \emph{negative}. The model captures how large TLE error is in
aggregate but gets the per-event ranking backwards. The honest claim is
that it reproduces the scale of uncertainty, not which conjunctions are
worst. Cause: a single scalar observable per event is dominated by the
reported miss magnitude, and the model's per-event variation keys off the
same quantity in the opposite direction.}

\failed{Residual-versus-TLE-age correlation is only $-0.16$, so the data
does \emph{not} support an explicit time-growth term. Staying with the
fixed-scale model is an empirical decision, not a simplification.}

\subsection{Monte Carlo and its cost}

The short-term encounter model applies: at
\SIrange{10}{15}{\kilo\meter\per\second} relative speed the encounter
lasts milliseconds, relative motion is effectively rectilinear, and the
three-dimensional problem collapses onto the two-dimensional encounter
plane perpendicular to the relative velocity. $\Pc$ is then the integral
of the projected Gaussian over a disk of the combined hard-body radius.

Three independent estimators were implemented: batched Monte Carlo, polar
quadrature, and a small-disk closed form valid when
$\text{HBR} \ll \sigma$.

\worked{Monte Carlo and quadrature agree within Monte Carlo error bars on
every real event, maximum $|z| = 1.4$. Two methods sharing no code path,
so the encounter geometry and the sampler are cross-validated.}

\begin{table}[htbp]
\centering
\caption{Brute-force Monte Carlo convergence at $\Pc = 4.5\times10^{-6}$.
         This is the cost the surrogate replaces.}
\label{tab:mccost}
\begin{tabular}{rlS[table-format=3.1]}
\toprule
{Draws} & {$\Pc$ estimate} & {Relative error (\si{\percent})} \\
\midrule
\num{10000}     & \num{0.0}          & 100.0 \\
\num{100000}    & \num{0.0}          & 100.0 \\
\num{1000000}   & \num{6.0e-6}       & 33.1 \\
\num{10000000}  & \num{2.7e-6}       & 40.1 \\
\num{50000000}  & \num{4.56e-6}      & 1.2 \\
\bottomrule
\end{tabular}
\end{table}

At \num{100000} draws the estimator returns \textbf{zero hits} -- not a
poor answer, but no answer. Reaching \SI{10}{\percent} relative error
requires $\sim$\num{22000000} draws and \SI{1}{\percent} requires
$\sim2.2\times10^{9}$, per conjunction, against a catalogue of thousands
of screened pairs.

\subsection{Known caveats}\label{sec:caveats}

\begin{enumerate}
  \item Computed $\Pc$ runs $10$--$100\times$ \emph{below} the published
        \texttt{PC}. The dominant driver is almost certainly the hard-body
        radius: RCS-derived values of \SIrange{1}{4}{\meter} are likely
        below operational values, and $\Pc \propto \text{HBR}^2$, so a
        $10\times$ error in HBR is $100\times$ in $\Pc$. Agreement with
        \texttt{PC} must \emph{not} be presented as validation.
  \item Per-event calibration correlation is negative (above).
  \item One event showed a \SI{20.8}{\kilo\meter} rebuild error, far
        outside plausible TLE error. Likely a manoeuvre or stale elements;
        should be identified before it contaminates Phase 4.
  \item Slow encounters (relative velocity below
        \SI{1}{\kilo\meter\per\second}) violate the short-term encounter
        assumption and are currently filtered out.
\end{enumerate}

% =====================================================================
\section{Phase 3: space-filling designs and the surrogate}\label{sec:phase3}

\subsection{A negative result that redirected the phase}\label{sec:negative}

The intuitive reading of ``use a space-filling design instead of Monte
Carlo'' does \textbf{not} work. $\Pc$ is a rare-event indicator integral
over a microscopic disk. A 256-point design returns zero hits just as
\num{100000} random draws do, and in fact more reliably, because a
space-filling design deliberately spreads points \emph{away} from any
small region. Quasi-Monte Carlo does not rescue this either: its
advantage requires a smooth integrand, and an indicator function over a
tiny region is the opposite.

The construction that does work is a surrogate over the
\textbf{encounter-parameter space}, trained on expensive evaluations at
design points and amortised across the catalogue. This is recorded because
it is the kind of dead end worth reporting in a write-up rather than
quietly discarding.

\subsection{Dimensional reduction, verified rather than assumed}

$\Pc$ depends on six raw parameters: the miss vector $\mu$ (2), the
projected covariance $C$ (3), and the hard-body radius $R$ (1). Only
\textbf{four} are independent, because the defining integral
\begin{equation}
  \Pc = \int_{\|x\| \le R} \mathcal{N}(x;\, \mu,\, C)\, \mathrm{d}x
  \label{eq:pcdef}
\end{equation}
has a rotation-symmetric domain and is covariant under uniform scaling of
all lengths. Rotating into the eigenbasis of $C$ removes one parameter;
dividing every length by $R$ removes another, leaving
$(\mu_1/R,\, \mu_2/R,\, \sigma_1/R,\, \sigma_2/R)$.

This is an exact identity, and it was verified numerically rather than
argued from algebra (Table~\ref{tab:invariance}).

\begin{table}[htbp]
\centering
\caption{Numerical verification of the invariances underlying the
         $6 \to 4$ reduction. Errors are quadrature noise, not model
         error.}
\label{tab:invariance}
\begin{tabular}{lrS[table-format=1.1e-2]}
\toprule
Claim & {Cases} & {Max.\ relative error} \\
\midrule
$\Pc$ invariant under rotation of $\mu$ and $C$ & 300 & 8.3e-13 \\
$\Pc$ invariant under uniform length scaling    & 300 & 9.8e-13 \\
4-D reduction reproduces $\Pc$                  & 200 & 8.4e-13 \\
$\Pc$ even in each $\mu$ component (eigenbasis) & 200 & 6.2e-16 \\
$6 \to 4 \to \Pc$ round trip                    & 200 & 5.5e-13 \\
\bottomrule
\end{tabular}
\end{table}

\worked{The reduction is exact to machine precision, so it costs no
accuracy and saves substantial cost. Measured fill distance at $k=4$,
$n=64$ is \num{0.397}; matching that in six dimensions requires
$\sim$\numrange{500}{900} points, i.e.\ $8$--$14\times$ more expensive
evaluations for identical answers.}

The identity holds \emph{given the model}: circular hard body, Gaussian
uncertainty, short-term encounter. Slow encounters fall outside it. The
six-dimensional space is retained as an ablation so the saving can be
demonstrated empirically rather than asserted.

\subsection{The 4-D parameter space}

\begin{table}[htbp]
\centering
\caption{Primary parameter space. Ranges from 40 rebuilt conjunctions,
         widened. Anisotropy is parameterised as a ratio so
         $\sigma_2 \le \sigma_1$ holds by construction.}
\label{tab:paramspace}
\begin{tabular}{llll}
\toprule
Index & Parameter & Meaning & Range \\
\midrule
0 & $\log_{10}(\sigma_1/R)$        & uncertainty vs.\ object size & 32 -- 3162 \\
1 & $\log_{10}(\sigma_2/\sigma_1)$ & anisotropy                   & 0.032 -- 1 \\
2 & $\mu_1/\sigma_1$               & miss along major axis        & 0 -- 6 \\
3 & $\mu_2/\sigma_2$               & miss along minor axis        & 0 -- 6 \\
\bottomrule
\end{tabular}
\end{table}

\subsection{Design generation}

The RUSIS parallel-tempering MaxPro code was reused directly. The
pybind11 port compiled without modification and its designs come out
already scaled to $[0,1]$ as $(\ell - 0.5)/n$.

\worked{An independent \texttt{numpy} reimplementation of the MaxPro
criterion $\psi$ agrees with the C\texttt{++} extension to six decimal
places, cross-checking the original research code.}

\begin{table}[htbp]
\centering
\caption{Design quality at $k=6$. Lower $\psi$ and $\phi_p$ are better;
         larger minimum distance is better. Random LHD is the best of 20
         draws.}
\label{tab:designquality}
\begin{tabular}{rlSSS}
\toprule
{$n$} & Design & {$\psi$ (MaxPro)} & {$\phi_p$} & {Min.\ distance} \\
\midrule
64  & MaxPro (PT) & 37.7297 & 1.7649 & 0.3694 \\
64  & Random LHD  & 81.2452 & 2.3531 & 0.2684 \\
64  & Uniform     & 508.8451 & 2.8862 & 0.2118 \\
\midrule
256 & MaxPro (PT) & 79.6778 & 2.4753 & 0.2162 \\
256 & Random LHD  & 245.0499 & 3.4067 & 0.1533 \\
256 & Uniform     & 415.3693 & 3.3181 & 0.1615 \\
\bottomrule
\end{tabular}
\end{table}

Parallel tempering beats the best of 20 random Latin hypercubes by
$2.1$--$3.1\times$ on $\psi$, at \SIrange{0.2}{1.4}{\second} per design --
a cost paid once, since a design is a static artefact.

\subsection{Why training labels come from quadrature, not Monte Carlo}
\label{sec:labels}

$\Pc$ spans \textbf{12.8 orders of magnitude} across the parameter box
($2.1\times10^{-17}$ to $1.3\times10^{-4}$ on the $n=64$ design). Brute
force cannot label it:

\begin{table}[htbp]
\centering
\caption{Fraction of the 64 design points at which brute-force Monte Carlo
         yields a usable estimate (more than 10 expected hits).}
\label{tab:mcunusable}
\begin{tabular}{rr}
\toprule
{Draws per point} & {Points labelled (of 64)} \\
\midrule
$10^6$ & 6 \phantom{0}(\SI{9}{\percent}) \\
$10^7$ & 13 (\SI{20}{\percent}) \\
$10^8$ & 18 (\SI{28}{\percent}) \\
\bottomrule
\end{tabular}
\end{table}

Nineteen of 64 points lie below $\Pc = 10^{-12}$, which no Monte Carlo
will ever reach. Labels therefore come from the polar quadrature, which is
exact here and agrees with Monte Carlo wherever Monte Carlo works at all
(Section~\ref{sec:phase2}).

This strengthens rather than weakens the motivation: brute-force Monte
Carlo is not merely \emph{expensive} for rare events, it is
\textbf{unusable} across most of the parameter space. The surrogate
produces estimates where brute force produces nothing. The quadrature
serves as exact ground truth in this Gaussian testbed; in the general case
(non-Gaussian uncertainty, long-duration encounters, non-circular bodies)
no closed form exists and the expensive sampler is the only option.

\subsection{Surrogate}

Gaussian process with an anisotropic (ARD) squared-exponential kernel,
hyperparameters by maximising the log marginal likelihood. Targets are
$\log_{10}\Pc$ rather than $\Pc$: with 13 orders of magnitude of dynamic
range, absolute error is meaningless and a GP on the raw value would be
dominated by a handful of large-$\Pc$ points. The GP is plain
\texttt{numpy} (one Cholesky factorisation, two triangular solves); only
the optimiser comes from \texttt{scipy}.

% --- RESULTS TABLE INSERTED BELOW WHEN THE BENCHMARK RUN COMPLETES ---
% Included by project_log.tex. Regenerate the numbers with
%   ./venv/bin/python src/run_surrogate.py

\begin{table}[htbp]
\centering
\caption{Surrogate accuracy against a 3000-point held-out test set.
         RMSE is in $\log_{10}\Pc$, i.e.\ orders of magnitude. Every design
         type gets five replicates: parallel tempering is seeded from
         \texttt{random\_device}, so repeated runs give genuinely different
         designs. Spread is the standard deviation across replicates.}
\label{tab:surrogate}
\begin{tabular}{rlSSS}
\toprule
{$n$} & Design & {RMSE} & {Spread} & {Fit time (\si{\second})} \\
\midrule
32  & MaxPro (PT) & 0.034 & 0.001 & 0.15 \\
32  & Random LHD  & 0.046 & 0.024 & 0.18 \\
32  & Uniform     & 0.043 & 0.015 & 0.19 \\
\midrule
64  & MaxPro (PT) & 0.029 & 0.001 & 0.39 \\
64  & Random LHD  & 0.033 & 0.009 & 0.37 \\
64  & Uniform     & 0.045 & 0.013 & 0.28 \\
\midrule
128 & MaxPro (PT) & 0.018 & 0.003 & 0.89 \\
128 & Random LHD  & 0.022 & 0.004 & 1.06 \\
128 & Uniform     & 0.033 & 0.014 & 1.20 \\
\midrule
256 & MaxPro (PT) & 0.006 & 0.003 & 4.21 \\
256 & Random LHD  & 0.006 & 0.001 & 4.73 \\
256 & Uniform     & 0.005 & 0.001 & 5.32 \\
\bottomrule
\end{tabular}
\end{table}

\subsubsection*{Reading the benchmark honestly}

\worked{The surrogate is accurate in absolute terms. At $n=64$ the RMSE is
\num{0.029} orders of magnitude -- roughly \SI{7}{\percent} in $\Pc$ --
from 64 training evaluations, against a quantity spanning 13 orders of
magnitude. A single prediction costs \SI{8.3}{\micro\second} against
$\sim$\num{22000000} Monte Carlo draws for \SI{10}{\percent} accuracy at
one representative event.}

\worked{\textbf{The clearest advantage of parallel tempering is variance,
not mean accuracy.} At $n=32$ MaxPro's RMSE spread across replicates is
$\pm\num{0.001}$ against $\pm\num{0.024}$ for random Latin hypercubes --
more than twenty times more consistent. A random design is a lottery: some
draws are as good as an optimised one, others are much worse. With a
one-shot expensive evaluation budget, reproducibility is worth as much as
the mean improvement, and this is the practitioner-facing argument for the
method.}

\failed{\textbf{The $2$--$3\times$ advantage in the $\psi$ criterion does
not translate into a $2$--$3\times$ accuracy advantage.} Mean RMSE
improves by only $1.1$--$1.8\times$ over random LHD and uniform sampling
at $n \le 128$, and at $n=256$ the advantage disappears entirely
($0.99\times$ vs.\ random LHD, $0.74\times$ vs.\ uniform). Better
space-filling by the criterion is not the same as better surrogate
accuracy, and this should be stated plainly rather than glossed.}

\failed{At $n=256$ every design type converges to RMSE $\approx
\num{0.006}$. The target function $\log_{10}\Pc$ is smooth and
low-complexity, so once the box is adequately covered the design choice
stops mattering. \textbf{The interesting regime is small $n$}, which is
also the regime that matters when each evaluation is expensive.}

\gotcha{An earlier version of this benchmark compared one parallel-tempering
design against \emph{one} random draw, and reported MaxPro losing at
$n=128$ ($0.94\times$). That was sampling noise in the baseline, not a real
effect. Comparing a deterministic method against a single random
realisation is not a valid comparison; all figures above use five
replicates per design type.}

\subsubsection*{Points needed to match MaxPro at $n=64$}

Taking MaxPro at $n=64$ (RMSE \num{0.029}) as the target, random Latin
hypercubes first reach it at $n=128$ and uniform sampling at $n=256$ --
nominally $2\times$ and $4\times$ the evaluation budget. Given the
replicate spreads in Table~\ref{tab:surrogate} and the coarse grid of
sizes tested, these factors should be read as indicative rather than
precise; a finer sweep of $n$ would be needed to state them with
confidence.


% =====================================================================
\section{Reproducing everything}\label{sec:repro}

\begin{longtable}{p{0.42\textwidth} p{0.5\textwidth}}
\toprule
\textbf{Command} & \textbf{Produces} \\
\midrule
\endhead
\texttt{python src/validate\_iss.py}    & ISS propagation sanity checks \\
\texttt{python src/probe\_cdm.py}       & \texttt{cdm\_public} field inventory \\
\texttt{python src/rebuild\_events.py}  & Table~\ref{tab:rebuild} \\
\texttt{python src/calibrate.py}        & Table~\ref{tab:calib} \\
\texttt{python src/run\_montecarlo.py}  & Table~\ref{tab:mccost} \\
\texttt{python src/run\_surrogate.py}   & Table~\ref{tab:surrogate} \\
\texttt{python tests/test\_paramspace.py} & Table~\ref{tab:invariance} \\
\bottomrule
\end{longtable}

Test suites: \texttt{tests/test\_tle.py} (31),
\texttt{tests/test\_covariance.py} (22),
\texttt{tests/test\_montecarlo.py} (18),
\texttt{tests/test\_paramspace.py} (15).

\end{document}
